\documentclass[11pt]{article}
\usepackage[margin=1.12in]{geometry}
\usepackage{amsmath,amssymb,amsthm,mathtools}
\usepackage{booktabs}
\usepackage[hidelinks]{hyperref}
\numberwithin{equation}{section}
\newtheorem{theorem}{Theorem}[section]
\newtheorem{proposition}[theorem]{Proposition}
\newtheorem{lemma}[theorem]{Lemma}
\newtheorem{corollary}[theorem]{Corollary}
\theoremstyle{definition}
\newtheorem{definition}[theorem]{Definition}
\newtheorem{assumption}[theorem]{Assumption}
\theoremstyle{remark}
\newtheorem{remark}[theorem]{Remark}
\newtheorem{example}[theorem]{Example}
\newcommand{\W}{W}
\newcommand{\Wp}{W^+}
\newcommand{\Wm}{W^-}
\newcommand{\Ric}{\operatorname{Ric}}
\newcommand{\Rm}{\operatorname{Rm}}
\newcommand{\Scal}{\operatorname{Scal}}
\newcommand{\vol}{\operatorname{vol}}
\newcommand{\tr}{\operatorname{tr}}
\newcommand{\Hess}{\operatorname{Hess}}
\newcommand{\E}{\mathcal E}
\newcommand{\Bach}{\mathfrak B}
\newcommand{\R}{\mathbb R}
\newcommand{\Sph}{\mathbb S}
\newcommand{\SO}{\operatorname{SO}}
\newcommand{\Id}{\operatorname{Id}}
\newcommand{\Yrd}{Y_{\mathrm{rd}}}
\renewcommand{\epsilon}{\varepsilon}
\title{Marked necks and chiral Weyl energy on smooth homotopy four-spheres}
\author{Edward B. Baker III}
\date{25 November 2026}

\begin{document}
\maketitle
\begin{abstract}
We give estimates for cutting a buffered spherical neck in a smooth
Riemannian homotopy four-sphere. The estimates control the Yamabe
constants of the capped pieces and the chiral Weyl cost under
critical first-order metric comparison. A signature obstruction
identifies the link of a productized quotient collar below a
one-chirality energy of \(9\pi^2\). In the positive Yamabe band
\(4\pi^2\le E_+\le8\pi^2-\delta\), cut-and-cap arguments transfer
a supplied smooth asymptotically Euclidean limit, or a suitably
marked finite-puncture base, to the original smooth type. We extract
a separating collar between widely separated chiral energy quantiles
for exactly Bach-flat sources. On a fixed physical minimizing
sequence, a sufficiently wide elliptic spherical chart with
vanishing full-curvature integral and one point at controlled
distance from the inner center determines both complete sides.
Fixed small distortion and harmonic coordinates on the central
annulus then suffice for positive caps and the minimizing
contradiction. For sources with small integrated
Bach norm, pointwise curvature can diverge on low-energy balls even
after Yamabe minimization. We give hypotheses for preserving a
supplied cut under perturbation, and an initial-metric comparison
outside a small exceptional set from local Ricci smoothing. We
also analyze the conformal and closed critical alternatives. The
results leave open the existence of an appropriate starting class
on every smooth homotopy four-sphere and the extraction of the
anchored chart and its critical gauge on a general minimizing
sequence.
\end{abstract}

\tableofcontents

\section{Normalization and recognition}

The parent treatment develops the chiral connection
formulation and annular Bach variations
\cite[Sections 11--15]{Baker2026TwistorSections}.
Physical volume and the recovery of a weakly marked conformal
base are studied in \cite{Baker2026MarkedLimits}. Here the
comparison geometry is tracked through a cap, including which
parts of the original manifold remain on either side of the cut.

All manifolds and metrics are smooth unless otherwise specified. A
homotopy four-sphere is closed, connected and in the smooth category.
The unit round curvature tensor has slots
\(R_{ijkl}=g_{ik}g_{jl}-g_{il}g_{jk}\).
On an oriented Riemannian four-manifold the Hodge star on two-forms
gives real rank-three bundles \(\Lambda^\pm\). We use the
Hilbert--Schmidt norms on the Weyl blocks and the four-index tensor
norm on \(\Rm\). Define
\begin{align*}
 E_\pm(g;U)&=\int_U|W^\pm|_{\rm HS}^2\,dV_g,
 & \Bach(g;U)&=\int_U|B_g|_g\,dV_g,\\
 E_\pm(g)&=E_\pm(g;M),
 & \E&=E_++E_- .
\end{align*}
Our Bach convention is
\(B_{ij}=\nabla^k\nabla^l W_{ikjl}+\tfrac12\Ric^{kl}W_{ikjl}\).
The integrated first variation for compactly supported \(v\) is
\cite[\S1]{BahuaudGuentherIsenbergMazzeo2025Bach}
\[
 \left.\frac{d}{ds}\right|_{s=0}\int |W_{g+sv}|_{\rm tens}^2dV_{g+sv}
   =-4\int\langle B_g,v\rangle_g\,dV_g .
\]
The positive round scalar curvature convention gives the diagonal
curvature blocks \(W^\pm+(R/12)\Id_{\Lambda^\pm}\).
In particular,
\begin{align}
 |W|_{\rm tens}^2&=4(|W^+|_{\rm HS}^2+|W^-|_{\rm HS}^2),\\
 8\pi^2\chi(M)&=E_+(g)+E_-(g)+\int_M (R^2/24-|\Ric_0|^2/2)\,dV_g,\\
 12\pi^2\tau(M)&=E_+(g)-E_-(g).
 \label{eq:signature}
\end{align}
With \(\Delta=\operatorname{div}\nabla\) the Yamabe form is
\(Q_g(u)=\int(6|du|_g^2+Ru^2)\,dV_g\) and the quotient has
denominator \(\|u\|_{4,g}^2\). Under \(\widetilde g=f^2g\),
\(Q_{\widetilde g}(u)=Q_g(fu)\) and
\(\|u\|_{4,\widetilde g}=\|fu\|_{4,g}\) for compactly supported
tests, by the conformal Laplacian formula \cite[\S2]{LeeParker1987Yamabe}.
The round value is \(\Yrd=8\sqrt{6}\pi\). Each chiral curvature
measure is conformally invariant since the Hodge star on two-forms
is unchanged and the Weyl endomorphism has weight \(-2\).
Varying \(f^2g\) by \(f^2v\) in the conformally invariant
four-index Weyl functional gives \(B_{f^2g}=f^{-2}B_g\) as a
covariant tensor. Thus \(\Bach(g;U)\) is conformally invariant.

We consider volume-one positive Yamabe minimizers \(g_i\) on
oriented homotopy four-spheres \(M_i\), with
\begin{equation}\label{eq:band}
 R_{g_i}=Y_i=Y([g_i]),\qquad
 4\pi^2\le E_i=E_+(g_i)=E_-(g_i)\le 8\pi^2-\delta,\qquad \delta>0.
\end{equation}
The term \emph{almost Bach-flat} means \(\Bach(g_i;M_i)\to0\).
We take subsequences with \(E_i\to e\) and \(Y_i\to y\) as needed.
Put \(D_i=\int |\Ric_{0,g_i}|^2dV_i\).
Gauss--Bonnet and the sharp Yamabe bound give
\begin{align}
 Y_i^2&=48(8\pi^2-E_i)+12D_i,\\
 &\sqrt{48\delta}\le Y_i\le\Yrd,\qquad D_i\le4E_i,\\
 \int|\Rm_{g_i}|^2_{\rm tens}dV_i
 &=8E_i+2D_i+Y_i^2/6=16E_i+Y_i^2/3-64\pi^2.
 \label{eq:budgets}
\end{align}
In particular a fixed number \(\Lambda\) bounds the full curvature
integrals. For a fixed oriented smooth type \(M\), let
\[
 I^+(M)=\inf_{g:\,Y([g])>0} E_+(g)
\]
over smooth metrics. The orientation does not change this infimum,
by \eqref{eq:signature}.

\begin{proposition}[Strict recognition and the closed critical gap]
\label{thm:closed-gap}
If \(Y([g])>0\) on a homotopy four-sphere and \(E_+(g)<4\pi^2\),
the smooth manifold is diffeomorphic to \(S^4\). If also \(B_g=0\),
either the class is conformally round or \(E_+(g)>4\pi^2\).
\end{proposition}
\begin{proof}
On a homotopy four-sphere the full tensor Weyl integral is \(8E_+\).
Theorem A of Chang, Gursky and Yang applies when this integral is
less than \(16\pi^2\chi=32\pi^2\) and the Yamabe constant is positive.
Its conclusions are \(S^4\) and \(\mathbb{RP}^4\). Simple connectivity
excludes the second \cite[Theorem A and Remark 2]{ChangGurskyYang2003}.
For a Bach-flat metric at equality, Theorem B of the same paper, on
a manifold of either other smooth type, yields a conformal model
given by \(\mathbb{CP}^2\) or a manifold isometrically covered by
\(S^3\times S^1\). The former has nonzero \(b_2\), and the latter
cannot cover a simply connected manifold. Hence at or below the
threshold the underlying manifold is standard, and the inclusive
Bach-flat sphere rigidity theorem gives \(W=0\)
\cite[Remark 2 and Theorem 2.1]{ChangChen2021RicciSphere}.
The simply connected closed locally conformally flat manifold
is conformally round \cite{Kuiper1949}.
\end{proof}

\section{Buffered caps and the spherical link}

A buffered collar is an oriented embedding of
\(( -\ell-2,\ell+2)\times L\), with \(\ell\ge1\) and round
\(L=S^3/\Gamma\) of sectional curvature one, for a finite free
\(\Gamma\le\SO(4)\). The product model is \(h_0=dt^2+\gamma_L\).
Each side has a central unit splice \(0<|t|<1\) and a retained
buffer \(1<|t|<1+\ell\). The middle hypersurface separates:
a connected complement would give a loop of intersection number
one with this two-sided hypersurface, contrary to \(H_1(M)=0\).
The connected side collars show there are exactly two components.

The fixed cap on the cover is \(b=ds^2+a(s)^2\gamma_{S^3}\).
To construct it, start with \(\widetilde a(0)=0\) and
\(\widetilde a'=\cos(s)\zeta(s)\), where \(\zeta\) is smooth,
nonincreasing, one near zero and zero near \(\pi/2\). Extend
\(\widetilde a'\) by zero. If its integral is \(A>0\), set
\(a(s)=A^{-1}\widetilde a(As)\). The center has a round germ
and the boundary has a unit product germ. The radial and
tangential sectional curvatures are \(-a''/a\) and
\((1-(a')^2)/a^2\). Consequently the scalar curvature is
positive, including at the center. Continue by a semi-infinite
product to obtain \(b_\infty\). The change of radius satisfying
\(d\log r=ds/a(s)\) expresses \(b_\infty\) as
\((a/r)^2(dr^2+r^2\gamma_{S^3})\) on \(\R^4\). The factor
extends positively at the origin by the round germ, where
\(r=c\tan(As/2)\). Thus \(W=B=0\) and the compact-test Yamabe
constant is \(\Yrd\). Indeed stereographic projection transfers
the round inequality and a cutoff of a round constant at distance
\(\rho\) from the missing point costs \(O(\rho^2)\) in Dirichlet
energy. These tests can be radial. The model divided by \(\Gamma\)
therefore has orbifold Yamabe constant \(\Yrd/\sqrt{N}\),
\(N=|\Gamma|\).

\begin{theorem}[One buffered cut]\label{thm:caps}
Assume \(Y([g])=Y_0>0\). For a spherical collar in a source
conformal representative \(h\), suppose \(\eta_2=\|h-h_0\|_{C^2(h_0)}\)
is small on the whole collar. The two capped metrics \(g_j\)
have, for either chirality,
\begin{align}
 E_\pm(g_j)&=E_\pm(g;U_j)+E_\pm(g_j;T_j),\\
 \sum_{j,\pm} E_\pm(g_j;T_j)&\le C_W\eta_2^2,\\
 Y([g_j])&\ge\min\{Y_0,(1-C_Y\eta_2)\Yrd\}- C_I\ell^{-3/2}.
 \label{eq:smoothcap}
\end{align}
Here \(U_j\) is the unchanged part of the retained side and
\(T_j\) its splice. With central \(C^4\) error \(\eta_4\le1\), the
sum of new splice Bach integrals is at most \(C_B\eta_4\).
The constants in these estimates are independent of \(\ell\).

A second version uses the Euclidean collar
\(A_\ell/\Gamma\), \(A_\ell=\{e^{-\ell-2}<|x|<e^{\ell+2}\}\),
\(t=\log|x|\). The lift of a source conformal metric is
\(m=\delta_{\rm Euc}+v\). Put
\begin{align*}
 a&=\|v\|_\infty\le 1/4,\quad d=\|Dv\|_{4,dx},\\
 K&=\|\Rm_m\|_{2,m;\,(T_1\cup T_2)/\Gamma},\qquad
 \Delta_W=C\big(aK+N^{-1/2}(d^2+d+a)\big).
\end{align*}
If \(q_\pm=E_\pm(g;T_1\cup T_2)\), the sum of new splice
energies in that chirality has square root at most
\(\sqrt{q_\pm}+\Delta_W\). Thus the net increase for the two
capped metrics combined, or for a product replacement on the
uncut closed source, is at most \(2\sqrt{q_\pm}\Delta_W+\Delta_W^2\).
These bounds also hold separately with the corresponding data of
one splice. For sufficiently small \(C_\ell(a+d)\),
\begin{equation}\label{eq:weakcap}
 Y([g_j])\ge
 \min\{Y_0,(1-C_\ell(a+d))\Yrd/\sqrt N\}
       -C_I N^{-1/2}\ell^{-3/2}.
\end{equation}
For \(N>1\) these capped spaces are orbifolds with the cap center
modeled on \(\R^4/\Gamma\). The curvature bound \(K\) belongs to \(m\).
\end{theorem}

\begin{proof}
Extend the collar gauge over the source by cutting off the logarithm
of its positive factor outside the buffers. In the smooth version
set \(h_{\rm sp}=h_0+\theta(|t|)(h-h_0)\) on each splice,
where the fixed smooth profile \(0\le\theta\le1\) is zero near
zero and one near one. This is a smooth positive metric and
matches the cap. The cap contributes zero Weyl and Bach
curvatures. The metric jets of orders two and four of
\(h_{\rm sp}\) differ from those of \(h_0\) by \(O(\eta_2)\)
and \(O(\eta_4)\), respectively. Weyl and Bach depend smoothly on
these respective jets. Bounded normalized splice volume and
conformal invariance give the claimed costs.

In Euclidean coordinates use \(n=\delta_{\rm Euc}+\theta v\)
before multiplying the collar by \(|x|^{-2}\). For the all-lowered
Riemann tensor the second-order part is constant-coefficient:
\[
 \Rm_{\mathrm{low}}(p)=L(D^2p)+Q(p)[Dp,Dp].
\]
This follows by differentiating the Christoffel symbols of the
first kind, with all inverse-metric derivatives contributing to
the quadratic term. Product differentiation, with bounded
derivatives of \(\theta\) on the fixed splice, gives
\[
 |\Rm_{\mathrm{low}}(n)-\theta\Rm_{\mathrm{low}}(m)|
 \le C(|Dv|^2+|Dv|+a).
\]
For a matrix \(p\) near the identity let \(P_p\) take an
all-lowered curvature tensor to its Weyl block in an oriented
\(p^{-1/2}\)-orthonormal frame and multiply it by
\((\det p)^{1/4}\). This linear map on curvature tensors
depends smoothly on \(p\), including both trace subtraction and
the Hodge projection. Write \(F(p)=P_p\Rm_{\mathrm{low}}(p)\).
Then
\[
 |F(n)-\theta F(m)|\le
 C(|Dv|^2+|Dv|+a+a|\Rm_{\mathrm{low}}(m)|).
\]
The squared Euclidean \(L^2\) norm of \(F\) equals the chiral
integral on the cover. Hölder and division by \(\sqrt N\) therefore
give \(\Delta_W\) in the quotient, and the triangle inequality gives
the square-root bound. Equivalently, with
\(d_{\rm down}=N^{-1/4}d\), the derivative terms in \(\Delta_W\)
are bounded by \(C(d_{\rm down}^2+N^{-1/4}d_{\rm down}+N^{-1/2}a)\).
Because multiplication by \(|x|^{-2}\) takes place after this
estimate, a curvature bound in a different nonconstant conformal
gauge would require separate justification.

For the local Yamabe test on a cap, transplant to the cover
\(b=b_\infty\), retaining the actual perturbation through
its buffer and cutting it off outside that buffer. Denote
the resulting comparison by \(k=b+e\). On the fixed collar,
\(\|e\|_\infty\le C a\) and
\(\|\nabla_b e\|_4\le C_\ell(d+a)\) in the Euclidean-input
version. The exact model satisfies
\(Q_b(w)\ge c\|w\|_{H^1(b)}^2\) and
\(Q_b(w)\ge\Yrd\|w\|_{4,b}^2\).
For \(\mu=dV_k/dV_b\), comparison of connections gives
\[
 \mu R_k-R_b=\operatorname{div}_b V+H,qquad
 |V|\le C|\nabla_b e|,\qquad |H|\le C(|e|+|\nabla_b e|^2).
\]
In fact \(C^a{}_{bc}=\Gamma(k)^a{}_{bc}-\Gamma(b)^a{}_{bc}\)
is a contraction of \(k^{-1}\nabla_b e\).
Contract \(\Ric(k)-\Ric(b)=\nabla_b C+C*C\) in \(\mu k^{-1}\).
Moving \(\nabla_b\) past that coefficient contributes only
quadratic first-derivative terms. The undifferentiated
contraction of the background Ricci tensor contributes \(O(|e|)\).
Integrating against \(w^2\) bounds the scalar-form error by
\[
 C\big(a\|w\|_2^2+\|\nabla_b e\|_4\|w\|_4\|dw\|_2
                  +\|\nabla_b e\|_4^2\|w\|_4^2\big).
\]
Gradient and denominator comparisons have relative error \(O(a)\).
Consequently \(Q_k(w)\ge(1-C_\ell(a+d))\Yrd\|w\|_{4,k}^2\)
on the cover. For invariant tests the form divides by \(N\)
and the squared \(L^4\) norm by \(\sqrt N\).
In the \(C^2\) version one compares scalar and gradient
densities pointwise with \(b_\infty\), whose scalar curvature
has a positive lower bound. This gives the coefficient
\((1-C_Y\eta_2)\Yrd\) uniformly in the buffer length.

For either local comparison choose a cutoff angle
\(\alpha((|t|-1)/\ell)\) going smoothly from zero to \(\pi/2\)
in the unit angle interval, and use its sine \(\chi\) and
cosine \(\xi\). Thus \(\chi u\) is extended by zero to
the original closed source and \(\xi u\) is a test on the cap
comparison. The identity
\[
 Q_{g_j}(u)=Q_{g_j}(\chi u)+Q_{g_j}(\xi u)
 -6\int (|d\chi|^2+|d\xi|^2)u^2\,dV_{g_j}
\]
has error at most \(C_I N^{-1/2}\ell^{-3/2}\|u\|_4^2\):
the derivative squares are \(O(\ell^{-2})\) on volume
\(O(\ell/N)\). Also
\(\|u\|_4^2\le\|\chi u\|_4^2+\|\xi u\|_4^2\). This proves
the Yamabe bounds. For these fixed cutoff profiles the power
\(\ell^{-3/2}\) in the error operator is attained on a long exact
product by taking \(u=\varphi((|t|-1)/\ell)\) supported where the
angle derivative is bounded away from zero.
\end{proof}

A useful sufficient variant is \(a\to0\) with bounded
\(L^2\) Cartesian second derivatives \(K_2=\|D^2v\|_2\) on
an annulus slightly larger than the splice and the retained
buffers. For a compactly supported cutoff \(\zeta\), integrate by
parts in \(\int\zeta^4|Dv|^4\) to get, with
\(Z=\|\zeta Dv\|_4\),
\[
 Z^4\le 4a\|D\zeta\|_4 Z^3+4aK_2 Z^2 .
\]
This gives \(Z^2\le C a(K_2+a\|D\zeta\|_4^2)\). On fixed
buffers it yields the first-order form estimate. On the fixed
unit splices the density error is bounded by
\(C[aK_2+a(K_2+a)+\sqrt{a(K_2+a)}+a]\) for \(N=1\).
The all-lowered formula yields the requisite bound on \(K\) in
this case. In both versions choose the buffer length first using
the lower Yamabe bound, then make the comparison sufficiently
small at that length. A Bach cost uses the separate fourth-order
estimate.

\begin{proposition}[A quotient obstruction on the closed source]\label{prop:link}
Let a smooth oriented integral homology four-sphere contain an
embedded oriented collar of a round space form \(S^3/\Gamma\).
If a metric on the same closed source is precisely the unit round
product near the slice and its global \(E_+\) is less than
\(9\pi^2\), then \(\Gamma=1\).
\end{proposition}

\begin{proof}
Let \(P_1,P_2\) be the sides. The link has \(H_2=0\) by
three-dimensional duality and finite fundamental group, so
Mayer--Vietoris gives \(H_2(P_j)=0\) and both signatures vanish.
Hillman's finite group case for a locally flat embedding in an
integral homology four-sphere allows \(1\), \(Q_8\), and the
binary icosahedral group of the Poincaré homology sphere
\cite[Theorem 5.2(1) and Corollary 4.4.2]{Hillman2024LocallyFlat}.
In the smooth category the last case is excluded by Rokhlin.
Indeed the sides are spin by restriction from the ambient
manifold. Gluing either to the negative \(E_8\) plumbing with
the matching (unique) boundary spin structure gives a closed
smooth spin manifold of signature \(\pm8\), incompatible
with divisibility by sixteen
\cite[Descriptions (1) and (4)]{KirbyScharlemann1979Poincare}%
\cite[\S1.2]{Rokhlin1971Submanifolds}.

For a free orthogonal realization of \(Q_8\) the central
involution is \(-\Id\). The skew-orthogonal quaternionic generators
\(I,J\) give an orthonormal basis \((v,Iv,Jv,IJv)\) from any unit
\(v\), identifying the representation up to orientation.
The odd signature eta of the antipodal quotient is zero.
Each of the three order-four cyclic subgroups has quotient eta
\(-1/2\) in the quaternion orientation by the cotangent formula
with the two generator angles \(\pi/2\)
\cite[Proposition (2.12)]{APS1975SpectralII}. Averaging equivariant
eta over group elements, using the common order-two intersection,
gives \(8\eta_{\rm sig}=3\cdot 4(-1/2)-2\cdot 2(0)=-6\).
Reversing the boundary orientation changes its sign.

The product-boundary signature theorem in the unhalved eta
convention on even forms gives
\[
 E_+(g;P_j)-E_-(g;P_j)=12\pi^2\eta_{\rm sig}(\partial P_j).
\]
It uses \(\tau(P_j)=0\) and the density in \eqref{eq:signature}
\cite[Theorem (4.14)]{APS1975SpectralI}. One of the two sides has
positive difference \(9\pi^2\), forcing at least this much
global \(E_+\), a contradiction.
\end{proof}

\section{Marked topology and transfer of energy}

\begin{theorem}[Finite spherical cuts]\label{thm:nonsplitting}
Cut a homotopy four-sphere \(M\) along finitely many disjoint
embedded three-spheres and fill every new boundary with
a marked smooth ball. Write \(m_j\) for the raw \(E_+\) of piece
\(j\), \(e_j\) for its capped energy, and assume the cap has
positive Yamabe constant and \(e_j\le m_j+\gamma_j\), \(\gamma_j\ge0\).
If \(E_+(g)+\sum_j\gamma_j<8\pi^2\), at most one capped piece
has \(e_j\ge4\pi^2\), and the remaining pieces are standard.
The original smooth type is that of the exceptional piece if it
exists and is standard otherwise. If \(M\) is nonstandard, the
exceptional piece \(s\) satisfies
\begin{equation}\label{eq:nonsplitting}
 \sum_{j\ne s}m_j\le E_+(g)-I^+(M)+\gamma_s.
\end{equation}
\end{theorem}

\begin{proof}
Each sphere separates, so the incidence graph for the
complementary pieces is a tree. Successively applying van Kampen
and Mayer--Vietoris along its edges, capping leaves fundamental
and second homology groups unchanged, and gives
\[
 \pi_1(M)=*_{j}\pi_1(\widehat M_j),\qquad
 H_2(M;\mathbb Z)=\bigoplus_j H_2(\widehat M_j;\mathbb Z).
\]
All summands are trivial. Poincaré duality and Hurewicz--Whitehead
therefore give homotopy spheres \(\widehat M_j\).
Proposition~\ref{thm:closed-gap} and the capped energy sum prove the
assertions on recognition and the number of exceptions.

Removing the actual caps and using the original boundary
identifications recovers \(M\). A parametrized embedded ball
in a recognized standard sphere can be shrunk by an
embedding isotopy to a small coordinate ball and straightened
there, extending the isotopy to the ambient manifold. The
component theorem for \(\operatorname{Diff}(S^3)\) allows every
attaching map to extend over the standard ball, by an isotopy
to an orthogonal map in the corresponding orientation component
\cite[p.~553 and the main theorem]{Hatcher1983Smale}. Using
standard leaves one at a time gives the oriented connected-sum
identification. The cap of the surviving piece is consequently
an admissible competitor on \(M\). Hence
\(I^+(M)\le e_s\le m_s+\gamma_s\), which proves \eqref{eq:nonsplitting}.
\end{proof}

For multiple caps on one retained piece the local Yamabe
estimates of Theorem~\ref{thm:caps} compose when the buffers
are disjoint. With cap coefficients \(y_b>0\) and lengths
\(\ell_b\), use one interior cutoff and one per cap, with a
sine/cosine pair on each buffer. The sum of their squares is
one, and the squared \(L^2\) norm of the combined gradient-square
density is at most \(C\sum_b\ell_b^{-3}\). Thus
\[
 Y([g_j])\ge \min\{Y([g]),\min_b y_b\}
                -C\Big(\sum_b\ell_b^{-3}\Big)^{1/2}.
\]
Each source or model term is tested after the corresponding zero
extension. If \(q_b\) is the raw splice energy and \(d_b\)
bounds its density error, the net excess per cap is at most
\(2\sqrt{q_b}d_b+d_b^2\). For disjoint splices their total is
at most \(2\sqrt{\sum q_b}\sqrt{\sum d_b^2}+\sum d_b^2\).
In particular at \(E_+(g_i)\to I^+(M)<8\pi^2\), positive caps
with vanishing excess force zero limiting energy on all raw
discarded pieces. If buffer length is also a parameter, first fix
it to obtain the Yamabe margins, take the sequence limsup at
that length, and then exhaust the buffers. For a single cut,
masses larger than \(E_+(g)-I^+(M)+\max_j\gamma_j\) on both
unchanged sides are impossible on a nonstandard \(M\).

\begin{theorem}[Finite ordinary punctures]\label{thm:finite-punctures}
Let \(g_i\) be smooth metrics on oriented homotopy
four-spheres \(M_i\), with \(Y([g_i])\ge y_0>0\) and
\(E_+(g_i)\to e\). Let \((N,q)\) be a smooth oriented homotopy
four-sphere with \(Y([q])>0\) and finitely many distinct marked
points. Delete disjoint normal balls of radius \(r\) and write
\(K_r\) for the resulting compact manifold with boundary.
Suppose that on a neighborhood of each fixed sufficiently small
\(K_r\) there are whole, oriented, smooth comparison embeddings
\(\Phi_{i,r}\) and positive factors taking \(\Phi_{i,r}^*g_i\)
to smooth convergence to \(q\). Put \(e_C=E_+(q)\). If
\(e-e_C<4\pi^2\), then \(N\) and \(M_i\) are orientation-preservingly
diffeomorphic for large \(i\), and \(I^+(M_i)\le e_C\le e\). Moreover
\begin{equation}\label{eq:exterior}
 \lim_{i\to\infty}E_+(g_i;M_i\setminus\Phi_{i,r}K_r)
   =e-e_C+E_+(q;N\setminus K_r)
\end{equation}
at each such \(r\). In particular a minimizing sequence has
vanishing exterior energy in the limit \(i\to\infty\) followed
by \(r\downarrow0\).
\end{theorem}
\begin{proof}
Let \(K\) be the embedded whole truncation. The injection
\(H_0(\partial K)\to H_0(K)\oplus H_0(M_i\setminus\mathring K)\)
from Mayer--Vietoris and connectedness of \(K\) show that each
exterior component has exactly one boundary sphere. The capped
exteriors are homotopy spheres by the tree argument.

Use normal radius \(\rho=r e^{-t}\) at each puncture, with a
fixed angular reflection to preserve orientation. In Cartesian
normal coordinates \(q-\delta_{\rm Euc}=O(\rho^2)\),
\(Dq=O(\rho)\), and higher derivatives are bounded. In product
coordinates \(\rho^{-2}q\) differs from the round product by
\(O_\ell(r^2)\) in \(C^2\): the factors from each derivative of
the coordinates have the required powers of \(\rho\). Take the
comparison on a larger fixed truncation including both buffers.
The source error in this gauge is then
\(\eta_{i,r}\le C_\ell r^2+o_i(1;r,\ell)\), on disjoint collars.
First fix \(\ell\) so the localization loss is less than
\(\tfrac12\min(y_0,\Yrd/2)\), then take \(r\) small and
\(i\) large. Theorem~\ref{thm:caps} gives positive exterior
Yamabe constants and total capped exterior energy at most
\[
 E_+(g_i)-E_+(g_i;\Phi_{i,r}K_r)+k C_W\eta_{i,r}^2
 < e-e_C+4\epsilon<4\pi^2 ,
\]
where \(k\) is the number of punctures and \(\epsilon>0\) is
chosen from the strict energy margin. Here the omitted smooth
tail costs less than \(\epsilon\), compact convergence and
global convergence each cost less than \(\epsilon\), and the
total splice cost is less than \(\epsilon\).

Each exterior cap is now recognized as standard. To match its
marking, double the deleted normal ball in \(N\). In the double
identify the actual cap ball in the recognized sphere with the
oppositely oriented copy, using an identification on a
neighborhood prescribed by the given normal collar. The
extension theorem for an embedded parametrized disk in the
standard sphere extends this local map on a neighborhood of
the disk to an ambient diffeomorphism
\cite[Theorem C]{Palais1960Extending}. Restrict to the exterior
and glue to \(\Phi_{i,r}^{-1}\) on the core. This proves the
claimed smooth identification with the actual puncture
markings. Pulling back \(q\) gives the admissible competitor.
Finally smooth convergence in the chosen conformal gauges and
invariance of chiral energy give \eqref{eq:exterior}. The
argument only requires fixed-radius comparisons. The ordered
exhaustion, or a diagonal chosen after those comparisons, needs
no rates at arbitrary moving radii.
\end{proof}

\section{A supplied smooth asymptotically Euclidean component}

\begin{theorem}[Same-type descent]\label{thm:ale-descent}
Assume \eqref{eq:band}, almost Bach-flatness, and a supplied
complete connected nonflat smooth pointed limit \((X,h)\)
at scales \(\lambda_i\downarrow0\). Suppose the comparisons
orientedly embed a neighborhood of every full compact
truncation, with smooth convergence of
\(\iota_i^*(\lambda_i^{-2}g_i)\). Assume finitely many
asymptotically locally Euclidean (ALE) ends, each of positive
order through five derivatives. Then there is one end with
trivial group. Its smooth one-point conformal compactification
\((\widehat X,\bar h)\) is, for all sufficiently large \(i\),
orientation-preservingly diffeomorphic to \(M_i\).
It is Bach-flat, has class Yamabe constant at least \(y\), and
\begin{equation}\label{eq:strictdrop}
 4\pi^2< e_X=E_+(\bar h)\le e-\kappa_\delta
 \qquad(\kappa_\delta>0).
\end{equation}
In particular no such limit occurs when \(M_i=M\) and
\(e-I^+(M)<\kappa_\delta\).
\end{theorem}

\begin{proof}[Proof of the end and compactification assertions]
Smooth convergence and scale invariance of integrated Bach
give \(R_h=B_h=0\). Transporting compactly supported Yamabe
tests shows that the compact-test constant \(Y_c(X,h)\) is at
least \(y\). Monotone exhaustion gives the budgets
\(e_\pm=E_\pm(h)\le e\) and
\(d=\int_X|\Ric_h|^2\le y^2/12-4(8\pi^2-e)\), the last by
using the source trace-free Ricci integral. In particular \(X\)
is noncompact, since scalar zero on a closed \(X\) would
contradict the constant test.

For \(k\) ends with groups of orders \(q_j\), the full core
\(K\) embeds in one homotopy sphere. Its boundaries are rational
homology spheres and have integral \(H_2=0\).
Mayer--Vietoris gives \(H_2(K)=0\), finite \(H_1(K)\), and a
distinct exterior for each end as before. The pair sequence
for \((K,\partial K)\), with \(H_3(K,\partial K)=H^1(K)=0\),
then gives \(H_3(K)=\mathbb Z^{k-1}\). Thus \(\chi(X)=2-k\)
and \(\tau(K)=0\). The Euler boundary contribution tends to
\(8\pi^2/q_j\) at each end, by rescaled \(C^2\) convergence
to the round Euclidean annulus and the curvature and second
fundamental form expression for the Euler boundary form.
The product-boundary signature formula follows by splicing
the cylindrical representatives to exact round quotient
products, with vanishing Weyl cost. These identities give
\begin{align*}
 e_++e_- -d/2 &= 8\pi^2(2-k-\textstyle\sum_j q_j^{-1}),\\
 e_+-e_- &= 12\pi^2\textstyle\sum_j\eta_{\rm sig}(S^3/\Gamma_j).
\end{align*}
The conventions are those of \cite[\S7, (7.2)]{TianViaclovsky2005ALE}
and \cite[Theorem (4.14)]{APS1975SpectralI}.
On the end of largest group order \(q_{\max}\), transplant
radial cutoffs of the Euclidean sharp test
\((1+|x|^2)^{-1}\) by dilation. First send the support out
the end at fixed cutoff radii, then remove the cutoffs, whose
Dirichlet errors tend to zero at zero and infinity.
The quotient divides by \(\sqrt{q_{\max}}\), so
\(y\le\Yrd/\sqrt{q_{\max}}\).
These estimates imply
\[
 (8\pi^2-e)+(e_++e_-)/2+4\pi^2(k-2+\textstyle\sum_jq_j^{-1})
 \le y^2/48\le 8\pi^2/q_{\max}.
\]
For \(k\ge2\) the leftmost side exceeds the rightmost by
at least \(\delta+4\pi^2(k-2)(1+q_{\max}^{-1})\).
Thus \(k=1\). The finite-group topology and spin argument in
Proposition~\ref{prop:link} apply to the embedded link. In the
\(Q_8\) case the end signature formula forces
\(|e_+-e_-|=9\pi^2\), outside the available budgets on both
blocks. Hence the end is AE and
\begin{equation}\label{eq:AEenergy}
 e_+=e_-=e_X,qquad d=4e_X .
\end{equation}

The compactified metric is scalar-flat near the added point.
On the AE end, a positive harmonic factor
\(a\) can be chosen with
\(D^j(a-r^{-2})=O(r^{-2-\beta-j})\) through order five for
some \(0<\beta<\min(1,\alpha)\), where \(\alpha\) is the
differentiated ALE order. One constructs a harmonic function
by solving on expanding annuli with boundary data \(r^{-2}\)
and trapping it between positive multiples of
\(r^{-2}\pm r^{-2-\beta}\). Indeed
\(\Delta_h r^{-2}=O(r^{-4-\alpha})\), and
\(\Delta_h r^{-2-\beta}=\beta(2+\beta)r^{-4-\beta}
+O(r^{-4-\beta-\alpha})\). Rescaled interior estimates
give derivative bounds for the trapped function. After cutting
it off near the inner end and extending it to \(\R^4\), its
Euclidean Laplacian is \(O(r^{-4-\alpha})\) at infinity.
Its Newtonian potential has a positive leading coefficient
times \(r^{-2}\), by the finite absolute \(\beta\)-moment of
the forcing and the lower barrier. To see the remainder
\(O(r^{-2-\beta})\), estimate the kernel difference on
\(|z|\le r/2\) by \(Cr^{-2-\beta}|z|^\beta\) and split the
remaining integral at \(|x-z|=r/4\) to use decay and tail
mass separately. The decaying entire harmonic difference
vanishes. Rescaled interior estimates for the difference
with its leading term give the differentiated error.
Normalize the coefficient to one, and extend \(a\) positively
over the core.

After inversion the scalar-flat metric \(a^2h\) differs from
Euclidean by \(O(\rho^\beta)\), with \(j\)th derivatives
\(O(\rho^{\beta-j})\) through order five. It is Bach-flat
by conformal covariance. The constant-scalar,
obstruction-flat removable point theorem, in dimension four
with trivial local group, supplies a smooth extending chart
\cite[Definition 1.6, Theorem 1.8 and Remark 7.2]{AcheViaclovsky2011Obstruction}.
It is this chart that is used for the smooth identification.
Capping the topological core shows simple connectivity and
integral homology of \(S^4\), by the separation and
van Kampen and Mayer--Vietoris argument. The compactification
is thus a smooth homotopy sphere. The point has zero
\(H^1\) cutoff cost, so conformal transport of compact tests
on the punctured manifold gives
\(Y([\bar h])=Y_c(X,h)\ge y>0\).
The chiral integrals are preserved.
\end{proof}

By \eqref{eq:AEenergy}, zero Weyl energy would give a flat
scalar-flat component, contrary to the hypothesis. The closed
gap gives \(e_X>4\pi^2\). The Ricci budget and
\eqref{eq:AEenergy} yield
\begin{equation}\label{eq:residual}
 e-e_X\ge (\Yrd^2-y^2)/48
         \ge (\Yrd^2-Y([\bar h])^2)/48>0.
\end{equation}
The strict inequality uses Schoen's equality theorem because
the compactified class has positive energy
\cite[Theorem 2]{Schoen1984Yamabe}. The upper bound
\(e-e_X<4\pi^2-\delta\) follows from the band and the closed gap.
The marked transfer theorem therefore applies with one ordinary
puncture, giving the diffeomorphism in Theorem~\ref{thm:ale-descent}.
Here the source has the required comparison at a fixed
normal radius \(s\) in \(\widehat X\): with
\(\rho=s e^{-t}\), the pullback is multiplied by
\((a/(\lambda_i\rho))^2\), whose normalized limiting metric on
the cylinder is \(\rho^{-2}\bar h\). The compact core and the
buffers lie in a single fixed embedded truncation. This proves
the qualitative descent without an energy identity on the
whole complement.

In the marked normal gauge the error through \(k\le5\)
product derivatives on \(|t|\le T\) is at most
\(C_{k,T}s^2+o_i(1;s,T)\), since
\(\rho^j|D^j(\bar h-\delta_{\rm Euc})|=O(\rho^2)\)
for \(0\le j\le5\) in the smooth normal chart. The
pre-transition subcore on the \(t=-1\) slice omits just
the normal ball of radius \(\exp(1)s\), so its incoming energy
tends to \(e_X\) as \(i\to\infty\) and then \(s\downarrow0\).
In the inverted end coordinates \(z=x/|x|^2\) the
comparison \( (1-C|z|^\beta)\delta_{\rm Euc}
\le \bar h\le (1+C|z|^\beta)\delta_{\rm Euc}\) likewise
gives \(\rho=|z|(1+O(|z|^\beta))\) by radial length
comparison. Since \(a=r^{-2}(1+O(r^{-\beta}))\) with
\(r=|x|\), the physical inverse conformal factor on a
slice is \((\lambda_i/\rho)(1+O_T(s^\beta))\).
Its \(g_i\) three-volume is therefore comparable to
\((\lambda_i s^{-1}e^t)^3\). This uses the inverted
coordinates for size and the smooth normal coordinates for
the extension of the marking.

\begin{lemma}[Exact compactness alternative]\label{thm:exact-compactness}
For exactly Bach-flat sources in \eqref{eq:band}, either a
subsequence converges smoothly via oriented diffeomorphisms
to a closed exact positive Yamabe minimizer or a nonflat
complete smooth ALE limit appears with the hypotheses of
Theorem~\ref{thm:ale-descent}. For every fixed band there is a
\(d_\delta>0\) such that all exact closed classes in the band
have \(Y\le\Yrd-d_\delta\).
\end{lemma}

\begin{proof}
In the Yamabe gauge, the scale-normalized Sobolev inequality is
\[
 \|f\|_4\le \sqrt{6/y_0}\|df\|_2+
                 \vol(M)^{-1/4}\|f\|_2,qquad y_0=\sqrt{48\delta}.
\]
Under constant rescaling keep the same coefficient of the
gradient term. Uniform small-ball noncollapse and a physical
diameter upper bound follow by the elementary argument in
Section~\ref{sec:quantiles}. Noncollapse passes to fixed radii
under increasing metric rescalings. If the curvature
maximum is bounded, work at the original scale. If it
diverges, rescale by the maximum and base at a maximum
point. The selected scales have uniform bounded curvature
and injectivity radius lower bounds by the local volume
bound and \cite[Theorem 4.7]{CheegerGromovTaylor1982}. In the
rescaled case their volumes diverge and diameter tends to
infinity by Bishop comparison. On sufficiently small
uniform balls the curvature energy is small, so exact
Bach-flat constant-scalar regularity bounds all curvature
derivatives \cite[Theorem 4.1 and the remark after Theorem 1.3]{TianViaclovsky2004Moduli}.
Smooth compactness gives embeddings of whole compact sets
in the pointed limit. In the bounded case the diameter
bound makes it closed, and a whole comparison embedding is then
a diffeomorphism. The limiting scalar of this unit-volume
metric is its Yamabe constant by passing the inequality to
fixed tests and taking the constant test. In the rescaled
case the limit has curvature norm one at the base point and
is scalar-flat. The complete pointed compactness theorem with
the exact normalized Sobolev coefficient, uniformly bounded
total critical curvature, lower total volume bound and
\(b_1=0\) gives finitely many ALE ends of order zero in the
Bach-flat setting \cite[Theorem 1.1, following remark and \S7]{TianViaclovsky2004Moduli}.
The scalar-flat obstruction-flat decay theorem improves these
to positive differentiated order
\cite[Definition 1.1 and Theorem 1.2]{AcheViaclovsky2011Obstruction}.

If a sequence of such exact metrics has \(Y_i\to\Yrd\),
a closed limit contradicts the closed Yamabe rigidity theorem
since its chiral energy is at least \(4\pi^2\). In the
rescaled case the inequality \(Y([\bar h])\ge y\) makes the
compactified class round. Equation~\eqref{eq:AEenergy} then
forces the rescaling to be flat, also a contradiction.
\end{proof}

For the quantitative part of Theorem~\ref{thm:ale-descent},
apply Lemma~\ref{thm:exact-compactness} to the compactified
class itself, which is exactly Bach-flat and in the band.
Equation~\eqref{eq:residual} gives
\[
 \kappa_\delta=\big(\Yrd^2-(\Yrd-d_\delta)^2\big)/48>0 .
\]
It is a non-effective constant. If the exact class is empty,
a nonflat component of the stipulated type is already excluded.

\begin{proposition}[An expanding portion of one supplied end]
Assume a supplied smooth AE comparison as in
Theorem~\ref{thm:ale-descent}, with end order \(\alpha>0\).
There is a single sequence of whole-core embeddings and
radii \(R_i\to\infty\) with \(\lambda_i R_i\to0\) such that
\[
 \limsup_{i\to\infty}
 E_\pm(g_i;\iota_i\{A\le |x|\le R_i\})
 \le C A^{-2\alpha}
\]
for every large fixed \(A\). The corresponding full Weyl
bound holds as well, and at fixed \(A\) the chiral integrals
converge to the energy of the end beyond \(A\).
\end{proposition}
\begin{proof}
For each \(R_m=e^m\) choose a comparison on the full core
through \(eR_m\). Above a threshold in \(i\), require
\(\lambda_i R_m\le R_m^{-1}\) and comparison error at most
\(R_m^{-\alpha}\) in the weighted norm
\(\max_{j\le4}\sup r^j|D^j(\iota_{i,m}^*(\lambda_i^{-2}g_i)-h)|\)
on that buffered annulus. Take increasing thresholds and
select the largest eligible \(m\) at each \(i\).
For \(r=e^t\), Cartesian errors in the cylindrical metric
\((\lambda_i r)^{-2}\iota_i^*g_i\) have the form
\(v_{ab}(r\omega)(\omega^a dt+d\omega^a)
(\omega^b dt+d\omega^b)\). Each product derivative through
order four is bounded by a sum of the weighted Cartesian
derivatives, so the cylindrical error is \(O(r^{-\alpha})\)
uniformly up to the outer buffer. The conformally flat product
has fixed slice volume, and integration of
\(C r^{-2\alpha}d\log r\) gives the bound. Fixed smaller
annuli have smooth convergence and the remaining tail obeys
the same bound, proving convergence at fixed \(A\).
The diagonal imposes no relation between \(R_i\) and any
separately supplied outer component.
\end{proof}

\section{Two quantiles and a separating comparison}\label{sec:quantiles}

Put \(y_0=\sqrt{48\delta}\).
For \(V(r)=\vol_{g_i}B(x,r)\) the Yamabe test equal to one
on \(B(x,r/2)\), zero past \(B(x,r)\), and with gradient at
most \(2/r\) gives \(y_0 V(r/2)^{1/2}\le (24/r^2+\Yrd)V(r)\).
For \(v(r)=V(r)/r^4\) and \(r\le r_0\) this reads
\(v(r)\ge c v(r/2)^{1/2}\), with
\(c=y_0/[4(24+\Yrd r_0^2)]\). Iteration and the Euclidean
small-ball asymptotic at each individual center give
\(V(r)\ge c^2 r^4\) uniformly. On a minimizing diameter
geodesic of length \(\operatorname{diam}(M_i,g_i)\), mark off points
spaced \(2\rho\), \(0<\rho\le r_0\). There are
\(1+\lfloor\operatorname{diam}(M_i,g_i)/(2\rho)\rfloor\)
disjoint open \(\rho\)-balls. Unit total volume yields
\(\operatorname{diam}(M_i,g_i)<2/(c^2\rho^3)\).

For a center \(p_i\), distance spheres have four-volume
zero: the distance function has unit gradient almost
everywhere off \(p_i\), whereas the gradient of a Lipschitz
function vanishes almost everywhere on any fixed level set.
The cumulative measure of balls for
\(\mu_i=|W^+|_{\rm HS}^2dV_i\) is therefore continuous.
Fix \(\eta>0\) with \(2\eta<E_i\) and choose quantile radii
\(a_i<b_i\) such that the inner ball and outer complement
each have mass \(\eta\). If \(b_i/a_i\to\infty\), let
\(N_i=\lfloor\log_2(b_i/a_i)\rfloor\). The middle half of
the intervening dyadic shells contains at least
\(\lfloor (N_i/2-2)/m_i\rfloor\) disjoint complete radial
blocks of \(m_i\) shells. Take a fixed \(m_i=m\) or
\(m_i=\lfloor\sqrt{N_i}\rfloor\). The curvature budget
\eqref{eq:budgets} selects a block with full curvature
integral at most \(\Lambda/\lfloor(N_i/2-2)/m_i\rfloor\).
Its endpoint radii \(\ell_i,u_i\) obey
\[
 a_i/\ell_i\to0,\qquad u_i/b_i\to0,qquad u_i\to0 .
\]
The last follows from the physical diameter bound.
The geometric-mean rescaling in the growing case, by
\(\rho_i^{-2}\) with \(\rho_i=\sqrt{\ell_i u_i}\), makes
the selected block cover every fixed positive rescaled radial
interval. All of these conclusions follow from the physical
Yamabe and curvature bounds.

\begin{proposition}[A sufficient weak comparison]\label{prop:weakquantiles}
Suppose a fixed-width embedded oriented cone chart in the
selected block has its middle link separating the two
complete quantile sets. On a fixed annulus, including
both buffers, assume its constant-scale physical metric
lift is \(\delta_{\rm Euc}+v_i\), with
\(\|v_i\|_\infty+\|Dv_i\|_4\to0\).
Choose the width initially to allow the localization in
Theorem~\ref{thm:caps}. Then the link is spherical for large
\(i\), and one can cap both sides with positive Yamabe
constants and vanishing added chiral energy. In particular
such a chart is impossible on a minimizing sequence on a
fixed nonstandard \(M\). In the derivative hypothesis the
quotient norm can replace the lifted norm.
\end{proposition}
\begin{proof}
The physical curvature budget bounds the splice \(K_i\)
and in the selected block gives \(K_i=o(1)\). Both
incoming splice energies are consequently \(o(1)\).
The transition estimate with downstairs norms makes the
middle an exact product on the closed source, at \(o(1)\)
cost. Proposition~\ref{prop:link} then forces trivial link
group. At that stage the two derivative norms coincide.
Use the Yamabe estimate, with the buffer length fixed
before the comparison errors tend to zero. In
Theorem~\ref{thm:nonsplitting} each raw side has at least
\(\eta\) because the chart and splices are between the
two mass sets. This contradicts minimizing energy.
\end{proof}

\begin{theorem}[Exact separating extraction]\label{thm:exactcut}
If \(B_{g_i}=0\) in \eqref{eq:band} and \(b_i/a_i\to\infty\)
for the pair of quantiles above, then, for any fixed
logarithmic half-width, a subsequence has an oriented
embedded spherical collar inside the growing selected block.
After a constant physical rescaling its lift converges
smoothly to the Euclidean annulus, on a neighborhood of
the closed collar. The middle slice separates the complete
inner and outer quantile sets, which are outside the
modified strips. In particular separated quantiles of fixed
positive mass are impossible for an exactly Bach-flat
minimizing sequence on a fixed smooth homotopy sphere in
the band.
\end{theorem}

\begin{proof}
Use \(h_i=\rho_i^{-2}g_i\) at the midpoint scale of the
growing block. Its total volume tends to infinity. The
normalized Sobolev inequality in
Lemma~\ref{thm:exact-compactness}, the exact
constant-scalar Bach equation, uniform total curvature,
and \(b_1=0\) allow pointed compactness on this scale
\cite[Theorems 1.1--1.3, following remark and \S7]{TianViaclovsky2004Moduli}.
We use the broader Bach-flat version, with \(C^0\) orbifold
points and order-zero ALE ends at this stage. Its upper
ball-volume estimate in the unscaled gauge is uniform.
Exact local regularity on balls in the growing block gives
\[
 |\nabla^k\Rm_{g_i}|(x)\le C_k r^{-2-k}
        \left(\int_{\mathrm{block}}|\Rm|^2dV_i\right)^{1/2}
 \quad(r/2\le d_{g_i}(p_i,x)\le r)
\]
when \([r/4,2r]\) is inside the block, by applying the
estimate to a ball of radius \(r/8\)
\cite[Theorem 2.4]{ChangQingYang2006Gap}. In the
pointed rescaling this tends to zero on each fixed positive
radial interval. The ball lower bounds and curvature
control give local injectivity on those intervals
\cite[Theorem 4.7]{CheegerGromovTaylor1982}.
Thus the pointed limit is flat and smooth off its marked
point, and noncompact since fixed large radial distances
occur in the source. The one-cone conclusion of pointed
compactness at possible singularities uses \(b_1=0\) and
the exact coefficient of the volume-normalized Sobolev
term. If the center is singular, the limit has finite
\(b_1\) by its compact core and finitely many ALE ends.
The Sobolev inequality on compact tests off the center
passes from the source, with the lower-order coefficient
tending to zero. The upper volume bound passes by smooth
exhaustion and pointed distance convergence. The local
punctured neighborhood is connected, and its curvature
integral vanishes. The constant-scalar Bach-flat
removability theorem thus gives a smooth orbifold metric
at the center \cite[Theorem 1.2 and \S5]{Streets2007Removal}.
It is flat there by continuity. At a smooth center use
ordinary normal coordinates. In either case a small
normal ball is Euclidean modulo a finite free spherical
group. Its radius is pointed distance on a smaller ball.

Fix a slightly larger annulus enclosing the desired
half-width in this normal ball. Smooth comparison on the
regular set of the pointed limit gives oriented embeddings
whose pulled-back metrics converge smoothly, with distances
to \(p_i\) converging uniformly to the model radius.
A further fixed constant rescaling sets the middle radius
to one. The image of a closed comparison cylinder with
model end radii \(r_-<r_+\) contains the \emph{whole}
source shell of radii \(\beta_-<\beta_+\) whenever
\(r_-<\beta_-<\beta_+<r_+\) are fixed. Otherwise an omitted
point sequence would converge in pointed distance to an
interior model point. That point's comparison images are
at distance tending to zero from the omissions, while the
distances to the comparison boundary stay bounded below
by radial distance separation. This contradicts any
path from the image to an omitted point. Every path from
the inner to the outer quantile set must cross the covered
shell. The middle link separates, and the complement of the
comparison cylinder has two connected components, one
adjoining each end. The connected inner ball is in one
component. The complete outer quantile set is in the
other, since a path within the first would have to meet
the covered shell. The comparison lies between the protected
radial endpoints in the growing block.

The small-cost product replacement in Theorem~\ref{thm:caps}
and Proposition~\ref{prop:link} identify the link as a
sphere. Choose the width large enough for positive Yamabe
caps and apply Proposition~\ref{prop:weakquantiles} for
the minimizing conclusion. The standard sphere has
\(I^+=0\), so the positive lower band already excludes
its minimizing sequence.
\end{proof}

The inner area hypothesis for the Chang--Qing--Yang
neck theorem follows here from the exact curvature estimate
and the uniform upper ball-volume ratio. At a protected
radius \(r\) each direction \(\theta\) minimizing to \(r\)
has \(\Ric(\dot\gamma,\dot\gamma)\ge-3k^2\) on \([r/2,r]\)
with \(k=C/r\). The cube root of the polar Jacobian,
\(y=J^{1/3}\), satisfies \(y''\le k^2 y\) before the
cut time. The Wronskian with
\(z(s)=\sinh(k(s-r/2))/k\) is nonincreasing from
\(-y(r/2)\), so \(y/z\) is nonincreasing.
Consequently \(J(r,\theta)\le C J(s,\theta)\) for
\(3r/4\le s\le7r/8\), including the endpoint by continuity.
Apply the area formula to all directions minimizing to
\(r\), and integrate the comparison before their cut
times. This gives
\(\mathcal H^3(S_r(p_i))\le C\vol B(p_i,r)/r\le C r^3\)
at every protected radius. In particular the window
\([4\ell_i,400\ell_i]\) and the trimming outer radius
\(u_i/4\) satisfy the area and scale premises of the theorem
for large \(i\), along with \(b_1=0\) and small annular
curvature. The proof above obtains the full smaller shell
using the pointed comparison and fixes the buffer after
taking that limit \cite[Theorem 3.1]{ChangQingYang2006Gap}.

\section{A physical shell located by one point}\label{sec:anchored-shell}

For the fixed-type minimizing problem take one smooth, closed,
oriented homotopy four-sphere \(M\), with \(g_i\) as in
\eqref{eq:band} and \(E_i\to I^+(M)\). Let \(r_i(\alpha)\) be
the least physical radius \(r\) for which the maximum of
\(\mu_i(B_{g_i}(x,r))\) over all \(x\in M\) reaches \(\alpha E_i\),
where \(\mu_i=|W^+_{g_i}|_{\rm HS}^2dV_i\). This maximum is
attained and continuous in \(r\). Indeed the ball masses are
jointly continuous in the center and radius by the zero mass of
distance spheres, and \(M\) is compact. On one subsequence with
\(a_i=r_i(\theta)\), \(b_i=r_i(1/2)\), and \(a_i/b_i\to0\), fix
\(0<\theta<1/2\) and take \(x_i\) attaining the smaller quantile.
Writing \(d_i=d_{g_i}(x_i,\cdot)\), the two full sets are
\begin{equation}\label{eq:anchored-marks}
 D_{i,-}=\{d_i<a_i\},\qquad D_{i,+}=\{d_i\ge b_i\},\qquad
 \mu_i(D_{i,-})=\theta E_i,\qquad \mu_i(D_{i,+})\ge E_i/2.
\end{equation}
The last estimate uses the maximum over centers at \(b_i\).
The block selection of Section~\ref{sec:quantiles} works with
these same sets, since it uses only the quotient of the radii
and the full physical curvature budget.

\subsection{Locating both complete sides}

Write \(A(u,v)=\{u<|z|<v\}\subset\R^4\), and let
\(\delta_{\rm Euc}\) be the Euclidean metric. The following
statement places the marks using a chart whose radial
position in the physical source need not be prescribed.

\begin{proposition}[Marking an anchored wide chart]\label{prop:anchored-marking}
Fix \(A_0<\infty\) and an ellipticity bound \(H\ge1\). There is a
fixed \(L>8\), depending only on \(A_0,H,y_0\), with the following
property. Suppose that \(a_i/\rho_i\to0\), \(\rho_i/b_i\to0\), and
that there are oriented smooth embeddings
\(\Phi_i:A(1/8,L)\to M\). Put \(S_i(r)=\Phi_i(\{|z|=r\})\) and
\(m_i=\rho_i^{-2}\Phi_i^*g_i\), and assume
\begin{equation}\label{eq:anchored-wide-data}
 d_i(S_i(1))\le A_0\rho_i,\qquad
 H^{-1}\delta_{\rm Euc}\le m_i\le H\delta_{\rm Euc},\qquad
 \int_{A(1/8,L)}|\Rm_{m_i}|^2\,dV_{m_i}\longrightarrow0.
\end{equation}
Here \(d_i(S)\) is the infimum of the physical distance over \(S\).
Then \(D_{i,-}\) and \(D_{i,+}\) are on opposite sides of \(S_i(1)\)
for all large \(i\), and both are disjoint from the entire closed
trim \(K_i=\Phi_i(\{1/4\le|z|\le4\})\).
\end{proposition}

\begin{proof}
Each embedded coordinate sphere has two connected sides. A
complementary path joining its two collar faces would produce
a loop crossing the sphere once modulo two, contrary to
\(H_1(M;\mathbb Z/2)=0\). Every complementary component abuts
the sphere and meets one of its two connected half collars.
The sides are labeled by their radial direction on the chart.
Similarly the image of an open radial band between two chart
spheres is a component of their complement. It is open by
the local inverse theorem and relatively closed there by
compactness of the closed band. A physical path crossing that
band from one face to the other has length at least
\(\rho_i H^{-1/2}\) times their difference in radii. This uses
only the part of the path between its last entry on one face
and its first exit on the other, so other parts may use ambient
shortcuts.

The diameter bound from Section~\ref{sec:quantiles} implies
\(\rho_i\to0\), and the ball lower bound there gives
\(\vol_{g_i}B(z,r)\ge\nu r^4\) for all \(0<r\le1\), for a uniform
\(\nu>0\) depending on \(y_0\). Choose a side \(U_i\) of \(S_i(1)\)
of physical volume at most one half. Define \(f_i=0\) on the
other side and let it rise linearly to one between coordinate
radii \(1\) and \(1/2\) or between \(1\) and \(2\), according to
the side chosen, extending by one on the rest of \(U_i\).
Ellipticity on this fixed strip gives
\(\int|df_i|_{g_i}^2dV_i\le C(H)\rho_i^2\). The Yamabe form,
applied by smooth density, gives
\[
 (1-2^{-1/2})y_0\|f_i\|_{4,g_i}^2
 \le 6\int|df_i|_{g_i}^2dV_i
\]
because \(\|f_i\|_{2,g_i}^2\le2^{-1/2}\|f_i\|_{4,g_i}^2\).
The plateau therefore has volume \(O(\rho_i^4)\), and the
strip has volume at most \(H^2|A(1/2,2)|\rho_i^4\). In all,
\(\vol_{g_i}(U_i)\le C_*(H,y_0)\rho_i^4\), with \(C_*\)
independent of the wider chart radius \(L\). For \(p\in U_i\),
the open ball of radius \(\min(1,d_{g_i}(p,S_i(1)))\) stays in
\(U_i\) by minimizing segments from \(p\). The lower volume
bound shows that every point of \(U_i\) is within
\((C_*/\nu)^{1/4}\rho_i\) of \(S_i(1)\) for large \(i\).

The coordinate-outer side contains the embedded shell
\(2<|z|<L/2\) of physical volume at least
\(H^{-2}\omega_4((L/2)^4-16)\rho_i^4\), where \(\omega_4\) is the
volume of the Euclidean unit ball. Choose \(L>8\) with
\begin{equation}\label{eq:anchored-width}
 H^{-2}\omega_4((L/2)^4-16)>C_*(H,y_0),\qquad
 (L/2-1)/\sqrt H>A_0 .
\end{equation}
The coordinate-inner side is then the smaller one. A nearest point
of \(S_i(1)\) to \(x_i\), unit-sphere chart paths of length at most
\(\pi\sqrt H\rho_i\), and the bound on the smaller side show that
every point of that side has physical distance \(O(\rho_i)\) from
\(x_i\). The whole trim has the same bound, using radial chart
paths as well. Thus \(D_{i,+}\) is entirely on the outer side,
outside the trim, even if it is disconnected.

It remains to locate \(x_i\). A point of the protected closed
band \(1/4\le |z|\le L/2\) has ambient distance at least
\(c(H)\rho_i\) from the complement of the chart. To exit, a
path must first cross one of the radial faces of the larger
compact band \([3/16,3L/4]\). If the inner ball met the
protected band, its entire radius-\(a_i\) ball would be in
the chart by this first-exit bound, since the distance between
any two of its points is at most \(2a_i=o(\rho_i)\). The
four-dimensional scale invariance of the curvature integral,
and \(4|W^+|_{\rm HS}^2\le |\Rm|_{\rm tens}^2\), would make
its mass tend to zero, contradicting \eqref{eq:anchored-marks}.
In particular \(x_i\) is outside that band. If it lay on the
coordinate-outer side, a minimizing path to \(S_i(1)\) would
cross the band from \(S_i(L/2)\) to \(S_i(1)\), at cost at least
\((L/2-1)\rho_i/\sqrt H\). This contradicts the anchor and
\eqref{eq:anchored-width}. The inner ball is connected by its
minimizing radii, so it belongs to the coordinate-inner side
and misses the trim. All cut locations have used the physical
metrics \(g_i\).
\end{proof}

There is a useful fixed-width alternative when the supplied trim
already lies between the two marks and within physical distance
\(t_i=o(b_i)\) of \(x_i\). The cutoff proof shows that no point
of the smaller side can belong to \(D_{i,+}\). Indeed the last
entry into the trim on a minimizing segment from \(x_i\) to the
unit sphere costs at least \(3\rho_i/(4\sqrt H)\) in remaining
radial length, so \(\rho_i=O(t_i)\). One oriented source arc
from \(x_i\) to a point of \(D_{i,+}\) with nonzero signed
intersection with \(S_i(1)\) now marks the entire median
exterior. The signed intersection records the side of the
endpoint relative to the starting point. For instance, if the
supremum of \(d_i\) on the inner trim face is strictly below
its infimum on the outer face, a minimizing segment to a median
exterior point that visits the open trim at an intermediate
distance has one outward crossing and no inward crossing.
Distance increases along that segment, so the visit must enter
through the inner face and exit through the outer one. In a
wide chart with the volume test in \eqref{eq:anchored-width}
and \(x_i\) outside the trim, the two infima
\(d_i(S_i(1/2))\le d_i(S_i(2))\) instead suffice to orient
the center. An outer center would have to cross the band
\([1/2,2]\) on the way to the nearer inner sphere, giving the
reverse strict inequality by the extra radial length.

\subsection{Critical coordinate control and the cap}

\begin{proposition}[A short critical shell cut]\label{prop:neck-critical-shell}
Suppose a chart on \(A(1/8,8)\) in the physical source already
has the complete marking of \eqref{eq:anchored-marks}, with both
marks off the closed trim \(K_i\). Assume its constant
rescaling \(m_i\) satisfies
\begin{equation}\label{eq:neck-critical-shell}
 \|m_i-\delta_{\rm Euc}\|_\infty\le\epsilon_*,\qquad
 \int_{A(1/8,8)}|\Rm_{m_i}|^2\,dV_{m_i}\to0,\qquad
 \|D T_i\|_{2,dz}\to0,\quad T_i^a=\Delta_{m_i}z^a .
\end{equation}
Here \(\epsilon_*>0\) is fixed sufficiently small on this annulus
before taking the source limit. In exact harmonic coordinates the
last condition is automatic. Such data are impossible on a fixed
nonstandard \(M\) at its positive-class infimum. Thus the wide
data in Proposition~\ref{prop:anchored-marking} together with
\eqref{eq:neck-critical-shell} on the central chart also give a
contradiction. The same implication holds with uniform and strong
\(W^{1,4}\) convergence to Euclidean in the supplied central chart
in place of the fixed small distortion and tension conditions.
\end{proposition}

\begin{proof}
Write \(v=m-\delta_{\rm Euc}\), \(\eta=\|v\|_\infty\) and
\(T^a=\Delta_m z^a\). The reduced metric equation has the form
\[
 m^{ab}\partial_{ab}v_{jk}
 =-2\Ric_{jk}-m_{ja}\partial_kT^a-m_{ka}\partial_jT^a
   +O(|Dv|^2),
\]
using \(T^a=-m^{bc}\Gamma^a_{bc}\)
\href{https://www.numdam.org/item/10.24033/asens.1405.pdf}{(DeTurck--Kazdan,
Lemma~4.1)}. For a compactly supported cutoff \(\chi\) set
\(X=\|\chi^2D^2v\|_2\), \(G=\|\chi Dv\|_4\). Integration by
parts against \(v\) in \(\int \chi^4|Dv|^4\) gives
\(G^2\le C\eta X+C_\chi\eta^2\). Apply the Euclidean
\(\|D^2f\|_2=\|\Delta f\|_2\) identity to \(f=\chi^2v\).
Freezing the principal coefficient costs \(C\eta X\), the
quadratic term costs \(CG^2\), and the commutators cost at most
\(C_\chi(G+\eta)\). Substituting the inequality for \(G^2\)
and absorbing with \(\eta\le\epsilon_*\) yields
\[
 X\le C_\chi(\|\chi^2\Ric\|_2+\|\chi^2 DT\|_2+\eta),\qquad
 G^2\le C_\chi\eta(\|\chi^2\Ric\|_2+\|\chi^2 DT\|_2+\eta).
\]
These Euclidean component norms are comparable to metric norms
for the indicated tensor sources. Use nested cutoffs protecting
the closed trim. Rellich compactness and these estimates give a
weak \(W^{2,2}\) limit \(h\), with strong \(W^{1,s}\) convergence
for \(s<4\) and almost-everywhere convergence. The quadratic
first-derivative terms converge in distributions by strong
\(W^{1,2}\). The bounded principal coefficients pass to the
limit against a fixed \(L^2\) test by dominated convergence.
The limit has zero right-hand curvature and differentiated
tension in the reduced equation, and zero full curvature in
distributions by the all-lowered curvature formula.

For interior regularity write the reduced equation on a ball as
\(h^{ab}\partial_{ab}h+C(h,Dh)Dh=0\), with drift coefficients
bounded by \(C_0|Dh|\). This is a small perturbation of the
Dirichlet Laplacian on \(W^{2,s}\cap W^{1,s}_0\), for \(s=2,3\).
The drift is small in \(L^4\) by the estimate above and is bounded
on that operator domain by the embedding into gradients in
\(L^{4s/(4-s)}\). A cutoff times \(h\) gives a forcing in
\(L^3\). Solving there and using uniqueness at exponent two
gives \(h\in W^{2,3}_{\rm loc}\). Its gradient lies in \(L^{12}\)
and the principal coefficients are continuous. Freezing them on
smaller balls gives \(W^{2,6}_{\rm loc}\), followed by ordinary
interior bootstrapping to a smooth flat metric. In the
approximately harmonic case a limiting constant coordinate
tension is compatible with this argument, which uses its derivative.
Subtract the reduced equations for \(m_i\) and \(h\), and localize
again. The principal perturbation and the drift of size
\(C(|Dm_i|+|Dh|)\) are absorbed at exponent two. The other terms
tend to zero in \(L^2\), including coefficient differences against
the smooth jets of \(h\), the curvature and tension sources, and
the cutoff commutators. For a drift commutator use the \(L^4\)
bound on the drift and strong \(L^4\) convergence of \(m_i-h\).
This gives strong \(W^{2,2}\) and hence strong \(W^{1,4}\)
convergence on smaller interiors.

The flat metric on the simply connected annulus \(A(1/4,4)\)
has an oriented developing map \(F\), obtained by integrating a
parallel orthonormal coframe. Since \((DF)^tDF=h\), the gradient
is uniformly within \(C\epsilon_*\) of \(\SO(4)\). The geometric
rigidity estimate for the fixed Lipschitz annulus gives a
rotation \(Q\) with \(\|DF-Q\|_2\le C\epsilon_*\)
\cite[Theorem 3.1]{FrieseckeJamesMueller2002Rigidity}.
After translation, Poincar\'e and the uniform derivative bound
make \(F-Qz\) uniformly small on the closed \(A(1/3,3)\) as
\(\epsilon_*\downarrow0\). In fact a fixed nonzero supremum
there would give a fixed lower \(L^2\) bound on an interior ball.
The straight homotopy gives degree one on \(A(1/3,3)\) for
each target point in \(A(1/2,2)\). The positive Jacobian makes
its preimage unique there. The pulled-back unit sphere separates
the original trim faces, since every path across the trim must
traverse the two intermediate radii \(1/3\) and \(3\), at which
the developed radii are below and above one, respectively. In
particular this change of coordinates retains both complete marks.

On the developed \(A(1/2,2)\) the metrics tend strongly to
Euclidean in \(W^{1,4}\) and are uniformly elliptic. The same
properties already hold on a central annulus in the alternative
first-order hypothesis. Apply the curvature comparison of
Theorem~\ref{thm:caps} with a fixed short splice cutoff, leaving
an exactly flat strip about the new cut. Uniform ellipticity and
the product rule for the all-lowered curvature tensor bound the
\(L^2\) error in the chiral densities by a constant times
\[
 \|\Rm_{m_i}\|_2+\|Dm_i\|_4^2+\|Dm_i\|_2+
 \|m_i-\delta_{\rm Euc}\|_2=o(1)
\]
in the developed chart. The projection change is bounded by the
first term, so uniform convergence after development is unnecessary.
The estimate holds along the interpolation on the closed source
and on each splice separately. A conformal multiplier on the
flat strip gives the exact round product at zero chiral cost.

The splice path remains positive. If a continuous smooth metric
path on the closed type beginning at a positive class first lost
that sign, the zero first eigenvalue of the conformal Laplacian
would give a scalar-flat representative by a positive
eigenfunction. Equation~\eqref{eq:signature} and Gauss--Bonnet
then require \(E_+\ge8\pi^2\), whereas the whole path here has
energy \(E_i+o(1)<8\pi^2\). Stretch the exact product to any
finite length on the same type by increasing its radial arclength
coefficient with a bump supported in the product. That path retains
the same chiral energy, so the Yamabe constants of the stretched
metrics are bounded below by a fixed positive constant, by
\eqref{eq:budgets} in their unit-volume Yamabe representatives.
For each fixed source index any compactly supported test on one
side with an attached half-cylinder extends by zero to a long enough
stretched source and inherits that inequality. Compactify the end
conformally using \(r=e^{-t}\), so \(r^2(dt^2+\gamma)\) is smooth
and flat at the added point. The vanishing point capacity in
dimension four transports the same Yamabe bound to the capped
side by approximation with compactly supported punctured tests.
The resulting caps are spherical and conformally flat. The
per-side nonnegative upper allowances \(\gamma_{i,\pm}\) above the
actual raw energies are \(o(1)\). Both raw masses are at least
\(\theta E_i\), and for all large \(i\)
\[
 E_i+\gamma_{i,+}+\gamma_{i,-}<8\pi^2,\qquad
 E_i-I^+(M)+\max\{\gamma_{i,+},\gamma_{i,-}\}<4\pi^2\theta .
\]
Theorem~\ref{thm:nonsplitting} contradicts these two raw masses.
All normalized windows and splice profiles are fixed in advance.
The auxiliary stretch for testing positivity is performed on an
exact product after the splice at each fixed index.
\end{proof}

\subsection{What is presently available from a physical block}

The physical Yamabe form also controls the number of exterior
components with a specified depth, if a transition-volume
bound is supplied. For \(R>2s>0\), suppose two components of
\(\{d_i>s\}\) reach depth \(R\). Such components correspond to
the through components in \(\{s<d_i<R\}\), by the
\(H_0\) Mayer--Vietoris sequence for the open cover by the
connected \(\{d_i<R\}\) and \(\{d_i>s\}\), using \(H_1(M)=0\),
and by minimizing segments to see that every exterior
component meets the annulus. On the smaller full exterior
component use the cutoff \(\min(1,(d_i-s)/s)\), extended by
zero. Its gradient energy is at most \(s^{-2}V_i(s)\) with
\(V_i(s)=\vol_{g_i}\{s<d_i<2s\}\). A ball about a point of
depth \(R\) of radius \(\min(1,R-2s)\) stays in the plateau.
The same half-volume Yamabe argument forces
\begin{equation}\label{eq:neck-physical-volume-test}
 V_i(s)\ \ge\ \frac{(1-2^{-1/2})y_0\sqrt\nu}{6}
                   s^2\min(1,R-2s)^2 .
\end{equation}
Supremal depth gives the same estimate by approximation. Thus
an upper bound \(V_i(s)\le C s^4\) rules out two such components
if \(R/s>2+\sqrt{C/c_Y}\) and \(R-2s\le1\), where \(c_Y\) is
the coefficient in \eqref{eq:neck-physical-volume-test}. The scalar-curvature
term of the full budget gives only
\(V_i(s)\le 6y_0^{-2}\int_{\{s<d_i<2s\}}|\Rm_{g_i}|^2dV_i\).
The block supplies an absolute \(o(1)\), without an \(O(s^4)\)
estimate as \(s\to0\). At fixed radii in the shrinking
normalized block, the needed physical rate is fourth order
in the scale. The positive-Ricci conformal gauge has its own
ball volumes and radii. Comparisons there require transfer of
both the conformal density and the distance levels before use
in \eqref{eq:neck-physical-volume-test}. Once an anchored full
chart is supplied, its smaller-side cutoff works directly.

On a fixed rescaled annular window \(h_i=\rho_i^{-2}g_i\), a
supplementary bounded \(L^p(h_i)\) norm of the physical Ricci
tensor for \(p>2\) gives local comparison. This norm in the
original metric is \(\rho_i^{2-4/p}\|\Ric_{g_i}\|_{L^p(g_i)}\)
on the corresponding physical window. Interpolation with
the vanishing critical curvature norm gives \(\Ric_{h_i}\to0\)
in \(L^q\), \(2<q<p\), on the same window. Apply integral
relative volume comparison on each protected ball with model
curvature zero to bound its upper volume ratio
\href{https://www.math.ucla.edu/~petersen/Papers/28.pdf}{(Petersen--Wei,
Lemma~2.3)}. The normalized Yamabe inequality then supplies
a uniform local homogeneous Dirichlet Sobolev inequality,
absorbing its \(Y_i\rho_i^2\|u\|_2^2\) term by the upper ball
volume. On smaller fixed balls the small critical curvature
and the scaled \(L^q\) Ricci control give harmonic coordinates
and \(W^{2,q}\) interior bounds by Yang's local theorem
\cite[Theorem 7.1]{Yang1992ConvergenceII}. Their interior limits
are flat by strong \(W^{1,4}\) convergence and the reduced equation.
These are protected local conclusions. Comparison of balls
centered at \(x_i\) would require Ricci input inside the omitted
core as well. Packing and gluing local charts into a wide
source annulus, identifying its link and anchoring it at the
same center remain separate requirements. A finite quotient
chart can enter Proposition~\ref{prop:link} under the stated
cost and product hypotheses. The anchored theorem concerns
the resulting spherical chart on the source.

\subsection{Logarithmic selection with physical Ricci comparison}
\label{subsec:physical-ricci-flux}

We give two Ricci comparison hypotheses under which the full
distance defect tends to zero on a selected physical block.
First, the following estimate converts a scalar volume error
and a one-sided Laplacian error into a bound on that defect.

\begin{lemma}[Finite radial defect estimate]\label{lem:physical-flux}
Let \(d\) be distance from a point in a smooth complete Riemannian
four-manifold with metric \(h\). Write \(V(t)=\vol_h\{d<t\}\) and
\(\mu=4dV_h-\Delta_h(d^2/2)\), with
\(\Delta=\operatorname{div}\nabla\). Fix \(0<a<a'<b'<b\) and
\(\eta\in C_c^\infty((a,b))\) with \(0\le\eta\le1\) and
\(\eta=1\) on \([a',b']\). For any constants \(C,c\) one has
\begin{equation}\label{eq:finite-physical-flux}
 |\mu|(\{a'\le d\le b'\})
 \le \|V-C-\omega_4c\,t^4\|_{L^\infty([a,b])}
     \int_a^b |5\eta'(t)+t\eta''(t)|\,dt
     +2\mu^-(\{a<d<b\}),
\end{equation}
where \(\omega_4\) is the Euclidean unit ball volume and
\(\mu^-\) is the negative part of \(\mu\).
\end{lemma}

\begin{proof}
The squared-distance Laplacian is a signed Radon measure
away from the center. Its singular part is nonpositive,
by the corresponding result for distance and the weak
chain rule
\cite[Corollaries~2.11--2.12]{MantegazzaMascellaniUraltsev2014DistanceHessian}.
The Lipschitz test \(\eta(d)\) admits smooth approximations
converging uniformly and strongly in \(W^{1,2}\).
Coarea and \(|dd|_h=1\) almost everywhere therefore give
\[
 \begin{aligned}
 \int\eta(d)\,d\mu
 &=\int_a^b(4\eta(t)+t\eta'(t))V'(t)\,dt\\
 &=-\int_a^b(5\eta'(t)+t\eta''(t))
          [V(t)-C-\omega_4c\,t^4]\,dt .
 \end{aligned}
\]
The constant drops out by integration by parts. The fourth
power contributes the integral of a multiple of \((t^4\eta)'\).
Use \(|\mu|=\mu+2\mu^-\) and \(0\le\eta\le1\) to obtain the estimate.
\end{proof}

\begin{proposition}[Two physical Ricci comparisons]
\label{prop:physical-ricci-flux}
Use the source and complete chiral marks in
\eqref{eq:anchored-marks}. Suppose one of the following
holds for a fixed \(K\ge0\).
\begin{enumerate}
\item \(\Ric_{g_i}\ge-3K g_i\) on \(B_{g_i}(x_i,b_i)\).
\item There are \(s_i\ge a_i\) with \(s_i/b_i\to0\), and
\(\Ric_{g_i}\ge-3K g_i\) on
\(\{s_i<d_{g_i}(x_i,\cdot)<b_i\}\).
For a uniform \(C_0<\infty\) suppose also
\[
 \vol_{g_i}\{2s_i<d_{g_i}(x_i,\cdot)<4s_i\}\le C_0s_i^4 .
\]
\end{enumerate}
There are scales \(\rho_i\to0\) with
\(a_i/\rho_i\to0\) and \(\rho_i/b_i\to0\) for which every
fixed compact strip of positive radii in
\(h_i=\rho_i^{-2}g_i\) has vanishing full-curvature integral
and vanishing total variation of
\[
 \mu_i=4dV_{h_i}-\Delta_{h_i}(d_{h_i}(x_i,\cdot)^2/2).
\]
If at these scales there is also a proper flat comparison
with the compact oriented source embeddings and radial
alignment of Proposition~\ref{prop:radial-distance-test},
then its whole spherical cut exists with both complete
marks off the closed trim. In particular the combined
hypotheses are impossible at the positive-class infimum of
a fixed nonstandard smooth type.
\end{proposition}

\begin{proof}
Let \(s_K(t)=\sinh(\sqrt K\,t)/\sqrt K\) for \(K>0\), with
\(s_0(t)=t\). The model volume is
\[
 v_K(t)=4\omega_4\int_0^t s_K(u)^3\,du .
\]
In the first case take the centered ratio
\[
 \Phi_i(t)=\frac{\vol_{g_i}B_{g_i}(x_i,t)}{v_K(t)} .
\]
Polar Riccati comparison along each minimizing ray makes
the polar Jacobian divided by \(s_K(t)^3\) nonincreasing
before its cut time, with initial value one. Assigning it
zero after the cut time and taking area and volume averages
gives a nonincreasing \(\Phi_i\le1\) on
\([a_i,\min(1,b_i)]\).

In the second case apply Riccati comparison on the part of
each uncut ray after any \(s_i+\varepsilon\). There the
model mean curvature \(3s_K'(t-s_i-\varepsilon)/
s_K(t-s_i-\varepsilon)\) starts with a pole. Let
\(\varepsilon\downarrow0\) and use the nonpositive singular
Laplacian to get the weak upper bound
\(3s_K'(d-s_i)/s_K(d-s_i)\) for \(\Delta_{g_i}d\) on the
protected annulus, with \(d=d_{g_i}(x_i,\cdot)\).
Testing radially and using coarea gives
\(DA\le 3(s_K'/s_K)(t-s_i)A\,dt\) for the coarea
area \(A(t)\). Thus \(A(t)/s_K(t-s_i)^3\) is nonincreasing.
Its weighted average from \(2s_i\) is
\[
 \Phi_i(t)=\frac{\vol_{g_i}\{2s_i<d<t\}}
                 {v_K(t-s_i)-v_K(s_i)} .
\]
The average decreases, and the measured transition volume
bounds it uniformly from above at \(4s_i\). Both cases have
a positive floor on the intervals used below. In the first
case this follows from the physical Yamabe small-ball
lower bound. In the second, for \(8s_i\le t\le\min(1,b_i)\),
a radius-\(t/8\) ball centered at distance \(t/2\) on a
minimizing segment reaching \(b_i\) lies in the numerator
annulus. The same small-ball bound applies.

In the first case let \(A_i=a_i\), in the second
\(A_i=8s_i\), and take \(B_i=\min(1,b_i)\) in both.
Then \(T_i=\log(B_i/A_i)\to\infty\).
Split the middle half of the logarithmic interval into
\(N_i=\lfloor\sqrt{T_i}\rfloor\) equal subintervals.
The curvature integrals add and the ratio drops telescope.
With \(\Lambda\) the source full-curvature bound, one
subinterval has the sum of its curvature integral and
its ratio drop at most
\[
 (\Lambda+\Phi_i(A_i)-\Phi_i(B_i))/N_i.
\]
Choose \(\rho_i\) at its exponential midpoint. Then
\(a_i/\rho_i,\rho_i/b_i\le e^{-T_i/4}\) and
\(\rho_i\to0\) by bounded diameter. The selected logarithmic
width tends to infinity, so the bound applies to every
fixed compact normalized strip. These steps for the
absolute errors only require the ratio upper bound.

Put \(d_i=d_{h_i}(x_i,\cdot)\) and \(V_i(t)=\vol_{h_i}\{d_i<t\}\).
In case one use \(\sigma_i=C_i=0\). In case two use
\(\sigma_i=s_i/\rho_i\to0\) and
\(C_i=\rho_i^{-4}\vol_{g_i}B_{g_i}(x_i,2s_i)\).
On every fixed strip \([a,b]\), model expansion after
rescaling gives
\[
 V_i(t)-C_i=\omega_4\Phi_i(\rho_i)t^4
        +O_{a,b}(N_i^{-1}+\sigma_i+K\rho_i^2).
\]
Indeed, in case one the model is \(v_{K\rho_i^2}(t)\).
In case two it is
\(v_{K\rho_i^2}(t-\sigma_i)-v_{K\rho_i^2}(\sigma_i)\).
The corresponding distributional upper bound for
\(\Delta_{h_i}(d_i^2/2)\) is \(dV_{h_i}\) times
\(1+3d_i s_{K\rho_i^2}'(d_i-\sigma_i)/
s_{K\rho_i^2}(d_i-\sigma_i)\). Its excess over four is
\(O_{a,b}(\sigma_i+K\rho_i^2)\). The normalized volume
of the strip is bounded by the volume profile with its
constant \(C_i\) removed. Lemma~\ref{lem:physical-flux}
now gives vanishing total variation on every inner trim.

Each smooth compactly supported test from the proper flat
comparison transports into one fixed trim and has bounded
supremum norm. The preceding bound gives
\eqref{eq:radial-distance-test} irrespective of the volume
omitted by the comparison image. One oriented comparison
embedding on a neighborhood of a compact buffered radial
strip containing the entire preimage of \([1/8,L]\)
gives the full link chart after radial rigidity.
Proposition~\ref{prop:radial-distance-test} then gives the
round quotient, excludes a nontrivial link group, and
places the full inner ball and full median exterior on
opposite sides of the cut. The first-order splice uses
the physical \(C^0\cap W^{1,4}\) convergence and the
physical curvature bound. Its two cap allowances are
\(o(1)\) in the same chirality as the source marks.
\end{proof}

Control of the scalar ratio alone gives a signed radial
average. The last term in \eqref{eq:finite-physical-flux}
is what permits control on all tests in the comparison
image. For the separate volume test in
Proposition~\ref{prop:radial-volume-test}, the two
normalized omissions must instead have vanishing
difference, with a positive retained image limit and
the same radial normalization on both measurements.
A fixed positive comparison endpoint requires the
shifted denominators from that proposition.

The Ricci bounds above concern the physical Yamabe
metrics. The positive-Ricci conformal representative has
different distances and volume densities. Its comparison
alone does not give the two physical hypotheses. In the
second case the measured transition rate \(O(s_i^4)\)
is stronger than the available absolute \(o(1)\) volume bound.
Both routes still require the proper compact source
embedding at the selected scale.

\subsection{Critical Ricci control with a polar Jacobian bound}
\label{subsec:critical-radial-defect}

On the pre-cut polar coordinates from the physical center \(x_i\),
write \(t=d_{g_i}(x_i,\cdot)\), \(l=\log t\), and
\(p_i(l,\vartheta)=t^{-3}J_i(t,\vartheta)\), where \(J_i\) is
the polar volume Jacobian. Let \(\Theta_{i,l}\) be the unit
directions whose cut time is strictly beyond \(e^l\), and
extend \(p_i\) by zero after each cut. The following condition
uses the same directions at two radii:
\[
 Q_i(s,l)=\int_{\Theta_{i,l}}
         \frac{p_i(l,\vartheta)^2}{p_i(s,\vartheta)}
         \,d\vartheta,\qquad s<l .
\]

\begin{proposition}[A critical radial comparison]
\label{prop:critical-radial-defect}
Use the physical sequence and complete marks of
\eqref{eq:anchored-marks}. Suppose there are intervals
\([A_i,B_i]\subset[\log a_i,\log b_i]\) of lengths
\(T_i=B_i-A_i\to\infty\) and fixed \(C<\infty\) and
\(0\le\gamma<1\) such that
\begin{equation}\label{eq:critical-polar-jacobian}
 Q_i(s,l)\le C e^{2\gamma(l-s)}
 \qquad(A_i\le s<l\le B_i).
\end{equation}
There are scales \(a_i\ll\rho_i\ll b_i\) on which the
physical full curvature integral and total variation of
\(4dV_{h_i}-\Delta_{h_i}(d_{h_i}(x_i,\cdot)^2/2)\)
vanish on every fixed compact strip of positive radii,
with \(h_i=\rho_i^{-2}g_i\).
If the proper flat comparison with compact first-order
embeddings of Proposition~\ref{prop:radial-distance-test}
is supplied at these scales, its round cut has the two
complete marks on opposite sides and off the closed trim.
In the minimizing band this yields the contradiction of
that proposition.
\end{proposition}

\begin{proof}
On a surviving ray set
\(w=t\Delta_g t-3\) and
\(f=t^2(\Ric_g(\partial_t,\partial_t))_-\).
Riccati and the trace inequality for the three-dimensional
shape operator give \((\log p)'=w\) and
\(w'+w+w^2/3\le f\) in log radius. The critical volume
identity on the pre-cut region is
\[
 f^2p\,dl\,d\vartheta=(\Ric_g(\partial_t,\partial_t))_-^2dV_g .
\]
Integrate \(w'+w\le f\) backwards from \(l\), then
integrate the result over all start times in \([s,l]\). This gives
\[
 w(l)_+\le
 \frac{\log^+(p(l)/p(s))}{e^{l-s}-1}
 +\int_s^l\frac{e^{u-s}-1}{e^{l-s}-1}f(u)\,du .
\]
All radii in this inequality are before the given ray's cut.
Put \(P(l)=\int p(l)\) and \(G(l)=\int w(l)_+p(l)\), with
angular integration on surviving directions, and
\(F(u)=\int f(u)^2p(u)\). Jensen on the final angular
measure \(p(l)d\vartheta/P(l)\) and Cauchy--Schwarz with the
earlier weight \(p(u)\) give
\[
 G(l)\le
 \frac{P(l)\log(1+Q(s,l)/P(l))}{e^{l-s}-1}
 +\int_s^l
   \frac{e^{u-s}-1}{e^{l-s}-1}\sqrt{F(u)Q(u,l)}\,du .
\]
The first term is taken to be zero if \(P(l)=0\).
By Fatou as \(s\uparrow l\), \eqref{eq:critical-polar-jacobian}
also gives \(P(l)\le C\). For \(l-s\ge1\) the first term is
bounded by \(C'(1+l-s)e^{-(l-s)}\), and the kernel multiplying
\(\sqrt{F(u)}\) is bounded by \(C'e^{-(1-\gamma)(l-u)}\).

Restore the index and take \(s=A_i\). On the middle half
\(J_i\) of the log interval, integrate this estimate. Since
\(\int_{A_i}^{B_i}F_i\) is bounded by a constant times the
physical full curvature budget, Fubini and Cauchy--Schwarz
give \(\int_{J_i}G_i=O(1+\sqrt{T_i})\).
The normalized coarea function \(P_i\) has distributional
derivative \(\int w p_i\,d\vartheta\,dl-d\kappa_i\), where
\(\kappa_i\ge0\) records the downward jumps of the polar
density at cut times. This follows by differentiating
along each ray, with zero extension after cut, and then
integrating in angle. Thus its positive variation is
at most \(\int G_i\), and its total variation on \(J_i\)
is at most \(C+2\int_{J_i}G_i\), since \(0\le P_i\le C\).

Divide \(J_i\) into \(\lfloor T_i^{3/4}\rfloor\) equal
open subintervals. On one, the sum of the full curvature
integral, \(\int G_i\), and the variation of \(P_i\)
is \(O(T_i^{-1/4})\). Let \(\rho_i\) be the exponential
of its midpoint. Its width in log radius tends to
infinity, and the middle-half buffers separate the
scale from both quantiles. Bounded diameter gives
\(\rho_i\to0\). Put \(l_i=\log\rho_i\).

For each fixed \(0<a<b\), the physical rescaled volume
profile \(V_i(t)\) on \([a,b]\), up to an additive
constant, is \(\tfrac14 P_i(l_i)t^4+o(1)\), uniformly.
Indeed \(V_i'(t)=t^3 P_i(l_i+\log t)\) almost everywhere
and one can take a one-sided representative of \(P_i\).
The singular part of the distance Laplacian is nonpositive
\cite[Corollaries~2.11--2.12]{MantegazzaMascellaniUraltsev2014DistanceHessian}.
Thus the negative mass of the full rescaled defect
on \((a,b)\) is just
\(\int_{\log a}^{\log b}e^{4u}G_i(l_i+u)\,du=o(1)\).
The finite estimate of Lemma~\ref{lem:physical-flux}
gives the total variation claim on each inner trim.
The comparison tests for Proposition~\ref{prop:radial-distance-test}
lie in such trims and have bounded supremum norms.
Consequently their transport tests tend to zero, including
when the comparison images omit physical volume. That
proposition gives the marked cut and cap comparison.
\end{proof}

The centered and protected Ricci hypotheses of
Proposition~\ref{prop:physical-ricci-flux} imply
\eqref{eq:critical-polar-jacobian} with \(\gamma=0\):
polar Riccati bounds the forward ratio on surviving rays,
and in the protected case the measured transition bounds
the shifted area ratio. The weighted condition can
also hold in the presence of pointwise Ricci spikes.
At the critical exponent the usual polar multiplier
\(\psi^3 J\), for \(\psi=(\Delta_g t-3/t)_+\),
has no coercive \(\psi^4J\) term. Its curvature term
therefore lacks the estimate used in supercritical
comparison \cite[Theorem~2.1]{DaiWeiZhang2018Local}.
Condition~\eqref{eq:critical-polar-jacobian} instead
controls the cost of transporting the critical integral
from an earlier radius. It and the compact comparison
embeddings remain additional to the physical ball lower
bound and full curvature budget.


\section{Almost Bach-flat sources}

\subsection{Curvature spikes in Yamabe gauge}

\begin{proposition}\label{prop:spikes}
Suppose the energies in an almost Bach-flat sequence
\eqref{eq:band} lie in the interior of the fixed band.
Fix any \(p>4\) and any \(q_i\) in any prescribed
nonempty open set \(O_i\) of the source. There are
perturbations \(m_i\) supported on arbitrarily small
coordinate balls about \(q_i\), with
\(\|m_i-g_i\|_{C^1}+\|m_i-g_i\|_{W^{2,p}}\to0\), using
the fixed source norms at each index. Every unit-volume
Yamabe minimizer \(\widehat g_i\) of \([m_i]\) obeys
the same band and almost Bach-flatness. Both curvature
measures \(E_\pm\) change by \(o(1)\) in total variation,
as do the Yamabe constants and the global full curvature
integrals in the physical gauge. There are radii \(\sigma_i\downarrow0\),
independent of the minimizing representative, with
\[
 B_{\widehat g_i}(q_i,2\sigma_i)\subset O_i,qquad
 \int_{B_{\widehat g_i}(q_i,2\sigma_i)}
      |\Rm_{\widehat g_i}|^2_{\rm tens}dV\to0,qquad
 \sigma_i^2|W_{\widehat g_i}(q_i)|_{\rm tens}\to\infty .
\]
\end{proposition}

\begin{proof}
On one fixed source \(g\), use oriented normal coordinates
about \(q\) and add the tensor
\[
 v_r=r^{2-\beta}\chi(x/r)(x_1/r)^2\,dx^2\otimes dx^2,
 \qquad 0<\beta<4/p ,
\]
with smooth \(\chi\) supported in the unit ball and one
near zero. Positivity holds for small \(r\).
The derivative bound \(|D^j v_r|\le C_jr^{2-\beta-j}\)
gives convergence in \(C^1\) and \(W^{2,p}\).
At \(q\) the zeroth and first jets vanish and the
only altered curvature component up to symmetries is the
\(12\)-sectional curvature increment \(k=\pm r^{-\beta}\).
Its \(1212\) Weyl increment is \(k/3\), since the added
Ricci eigenvalues are \(k,k,0,0\) and the scalar increment
is \(2k\). Thus \(|W_{g+v_r}(q)|\ge c_g r^{-\beta}\).
The altered curvature energy and total-variation Weyl
measure differences are \(O_g(r^{4-2\beta})\). Expanding
the local Bach expression in metric derivatives gives terms
of total derivative order at most four, with smooth
coefficients in the metric. Every term involving the change
is \(O_g(r^{2-\beta-4})\) on the support, so the integrated
Bach change is \(O_g(r^{2-\beta})\).

Here is the fixed-source normalization argument. More
generally let \(m_j\to g\) in \(C^1\cap W^{2,p}\), where
\(g\) is a volume-one positive Yamabe minimizer with
nonzero Weyl energy. Then \(Y=Y([g])<\Yrd\) by
\cite[Theorem 2]{Schoen1984Yamabe}. Scalar curvature
converges in \(L^p\), and the Yamabe forms differ by
\(o(1)(\|df\|_2^2+\|f\|_4^2)\) uniformly in \(f\) in the
fixed norms, with equivalent denominators tending to the
same norm. Since the limiting form controls these two
norms, the class constants converge by the infimum
inequality and the constant test. The unit normalized
minimizing factors \(u_j\) are bounded in \(H^1\) and
\(Q_g(u_j)\to Y\). Take a weak limit \(u\), strong in
\(L^2\) and almost everywhere on a subsequence, and put
\(a=\int u^4dV_g\). Fourth-power splitting
\cite[Theorem 1]{BrezisLieb1983Functional} and the
local sharp Sobolev inequality, with an arbitrarily small
loss in \(\Yrd\) and a lower order \(L^2\) term
\cite[Sections 2--4]{LeeParker1987Yamabe}, give
\[
 Y\ge Y\sqrt a+\Yrd\sqrt{1-a}.
\]
Indeed the gradient energy splits under weak convergence
and the \(L^2\) error tends to zero. Strict subroundness
implies \(a=1\), so convergence is strong in \(L^4\)
and \(H^1\).

For the equation
\(-6\Delta_{m_j}u_j=V_j u_j\) with
\(V_j=Y([m_j])u_j^2-R_{m_j}\), the potentials converge in
\(L^2\). At each fixed power \(q\ge4\), truncate \(V_j\)
so the \(L^2\) tail is uniformly small relative to
\(24(q-1)/q^2\). Testing by \(u_j^{q-1}\) and applying
the Sobolev inequality to \(u_j^{q/2}\), absorb the tail
and get \(\|u_j\|_{2q}^q\le C_q\|u_j\|_q^q\).
Iteration bounds the right-hand side in \(L^s\) for some
\(s>2\). The fixed \(g\) elliptic \(W^{2,s}\) estimate
passes to \(m_j\) by absorbing the operator difference,
small by \(C^1\) convergence. Thus \(u_j\) is bounded in
\(L^\infty\) and then in \(W^{2,p}\). Compactness in
\(C^1\) and the difference equation against the fixed
operator give strong \(W^{2,p}\) convergence to a smooth
strictly positive minimizing factor for \(g\).
This applies to any sequence of minimizing factors.
In particular at each fixed index all perturbed minimizing
factors have common upper, positive lower and \(W^{2,p}\)
bounds for sufficiently small perturbations, and their
distance in \(W^{2,p}\) to the compact set of fixed
source minimizers tends to zero.

Accordingly the perturbed Yamabe metrics at each fixed
index are uniformly bilipschitz to the source and their
full curvatures are bounded in \(L^p\), with constants
depending on the index. On an intrinsic ball of radius
\(2\sigma\), volume is at most \(C_i\sigma^4\) and the
curvature integral at most \(C_i'\sigma^{4-8/p}\).
First choose \(\sigma_i\) to place the ball in \(O_i\)
and make the integral small uniformly over these metrics.
Then choose \(r_i\) so small that
\(c_i\sigma_i^2r_i^{-\beta}\ge i\), which is possible
after conformal multiplication because the minimizing
factors have a uniform upper bound at the fixed index.
At the same time keep each energy within its strict band
margin and all changes of the stated norms, curvature
measures and Yamabe constants tending to zero. The chiral
measures and integrated Bach are conformally invariant,
and \eqref{eq:budgets} controls the global full curvature
integrals after Yamabe reselection.
\end{proof}

Proposition~\ref{prop:spikes} permits perturbations of an
almost-critical minimizing sequence and applies inside a
prescribed open part of a low-energy block. It rules out
the direct analogue of a pointwise scale-invariant
curvature estimate from small local full curvature and
small integrated Bach in that class. The perturbation is
small in critical first derivatives before reselection.
It gives no obstruction to a comparison at that regularity.

\subsection{Retaining a supplied cut after perturbation}

On an exact source as in Theorem~\ref{thm:exactcut}, let
\(\mathcal A_i\) be the transferred closed buffered
annulus, with constant physical comparison scale \(s_i\)
and complement \(U_i^-,U_i^+\). Its whole-shell property
gives constants \(0<c<C\) such that
\begin{equation}\label{eq:protected}
 B_{g_i}(p_i,c s_i)\subset U_i^-,\qquad
 \{d_{g_i}(p_i,\cdot)\ge C s_i\}\subset U_i^+,
 \qquad a_i/s_i\to0,\quad s_i/b_i\to0 .
\end{equation}
Indeed the small centered ball is connected and an
intermediate whole radial sphere in the collar prevents
its side from reaching farther than \(C s_i\) within the
complement. Consider a perturbation \(m_i\) and a
unit-volume minimizing metric \(\widehat g_i=u_i^2 m_i\).
Assume positive-band Yamabe bounds for the new source
and let \(\widehat \mu_i\) be its chiral measure.

\begin{proposition}[Quantile placement in the retained marking]
Assume centered distances satisfy
\(l_i d_{g_i}\le d_{\widehat g_i}\le L_i d_{g_i}\), with the
distances from \(p_i\) throughout. Suppose
\begin{equation}\label{eq:placement}
 L_i a_i<l_i c s_i,qquad L_i C s_i<l_i b_i,qquad
 \|\widehat\mu_i-\mu_i\|_{\rm TV}\to0 .
\end{equation}
Then for every fixed \(0<\eta'<\eta\), both complete new
quantile sets of mass \(\eta'\) lie in the respective
unchanged regions, and their radius ratio tends to infinity.
If the old sets lose no mass, the same holds at the level
\(\eta\) whenever \(2\eta<\widehat\mu_i(M)\).
Suppose additionally that in the original marked annulus
the rescaled metric \(k_i=s_i^{-2}\Phi_i^*m_i\) differs
from \(s_i^{-2}\Phi_i^*g_i\) by \(o(1)\) in
\(L^\infty\) and Cartesian first derivatives in \(L^4\),
and that \(\|\Rm_{k_i}\|_2=O(1)\) on the splices.
Then the old collar admits positive Yamabe caps with
\(o(1)\) Weyl cost after perturbation.
\end{proposition}
\begin{proof}
The new ball of radius \(L_i a_i\) contains the old
inner set yet is contained in \(U_i^-\) by the first
inequality and \eqref{eq:protected}. The new exterior
at radius \(l_i b_i\) similarly contains the old
outer set and is contained in \(U_i^+\). Both carry
mass at least \(\eta-o(1)\). By continuity select new
radii \(\widehat a_i\le L_i a_i\) and
\(\widehat b_i\ge l_i b_i\). The first placement
inequality gives \(\widehat b_i/\widehat a_i\ge
 b_i/(c s_i)\to\infty\). The argument at exactly \(\eta\)
is identical when neither mass decreases, for instance
when the perturbation is supported between the two sets.
Chiral measures are unchanged by Yamabe reselection.
The original splice measure was \(o(1)\), so the new
incoming splice measure is \(o(1)\) by total variation.
Apply Theorem~\ref{thm:caps} in the auxiliary metric
\(k_i\) before any nonconstant conformal gauge change.
The buffer length has been fixed with the positive
Yamabe margin. This gives both cap conclusions.
\end{proof}

A sufficient placement condition is
\(L_i/l_i=o(\min(s_i/a_i,b_i/s_i))\).
An additive alternative is
\(\sup_x|\lambda_i^{-1}d_{\widehat g_i}(p_i,x)-
d_{g_i}(p_i,x)|\le\epsilon_i s_i\) for some \(\lambda_i>0\).
Use new radii \(\lambda_i(a_i+\epsilon_i s_i)\) and
\(\lambda_i(b_i-\epsilon_i s_i)\), provided
\(a_i+2\epsilon_i s_i<c s_i\) and
\(C s_i+2\epsilon_i s_i<b_i\).
The same ratio tends to infinity. If the perturbations
are chosen by a fixed-source diagonal as in
Proposition~\ref{prop:spikes}, the new minimizing factor
can be made uniformly close to the compact set
\(\mathcal U_i\) of original unit-volume minimizing
factors. Length comparison gives
\(l_i=e^{-2\epsilon_i}\inf v_i\),
\(L_i=e^{2\epsilon_i}\sup v_i\) for some \(v_i\in\mathcal U_i\)
and \(\epsilon_i\downarrow0\).
Thus the multiplicative placement condition holds for
all reselected factors if
\(\sup_{v\in\mathcal U_i}\sup v/\inf v\)
obeys the displayed distortion bound, including if the
original minimizing metric is unique. For the localized
curvature spikes one may impose the annular comparison
at each fixed index by choosing the support still smaller,
even if it intersects a splice. The divergent Weyl jet
then allows arbitrarily large curvature in the normalized
collar chart while the cutting cost tends to zero.

\begin{example}[A change of the fixed marking]
Fix an exact balanced source as above. Move a point of
the old inner ball other than \(p_i\) to a point strictly
past distance \(b_i\) by an ambient isotopy supported in
a thin tube, avoiding \(p_i\) and a second point beyond
\(b_i\). The tube may have arbitrarily small Weyl measure
because an embedded arc has four-volume zero. Pulling
back by this orientation-preserving diffeomorphism leaves
the metric exact, the quantile radii about \(p_i\)
unchanged, and changes the measure in total variation
by at most twice the measure of the tube. The new
complete outer set has points in both original retained
regions, so the inherited slice fails to separate it.
Pulling back the collar along with the metric restores
separation. Intrinsic energy and Bach bounds alone
therefore give no control of distance margins in a fixed
marking, even with small changes in the curvature measure.
\end{example}

\subsection{Local smoothing on a protected block}

Use the geometric-mean rescaling on the growing physical
block in Section~\ref{sec:quantiles}. The lower ball-volume
ratios and the scale-normalized Sobolev inequality hold on
protected balls, and the full curvature norm is \(\epsilon_i\to0\)
there. To apply the local smoothing theorem of Lee and Lee,
assume also uniform \emph{upper} initial ball-volume ratios
at every center in the required protected neighborhood and
every radius through one, after a fixed constant dilation
if necessary. The proof uses the initial physical metrics
up to these constant rescalings.
On these normalized initial balls the theorem requires the
upper ratios, a uniform local inhomogeneous Sobolev
inequality and small critical full curvature, with no initial
Ricci lower bound \cite[Theorem 1.1]{LeeLee2026Local}.

\begin{proposition}[Initial metrics in smoothed coordinates]
Under the upper volume bound, there are oriented embeddings
of a fixed Euclidean ball into the
original sources obtained from time-\(t\) normal
coordinates. In these embeddings the pullbacks of the
initial metrics, with an additional factor \(t^{-1}\),
are within \(C(a+\epsilon_i+\sqrt t)\) of Euclidean
outside an exceptional set. Both the model volume of
that set and its volume in the rescaled initial metric
are bounded by \(C a^{-2}(\epsilon_i^2+t)\), with \(a\in(0,1]\).
For \(a=t^{1/4}\) these errors vanish by first taking
\(i\to\infty\) and then \(t\downarrow0\).
\end{proposition}
\begin{proof}
Write \(k_i(t)\) for the flow starting at the rescaled
physical source, with protected center \(q_i\). The
local smoothing theorem gives
\[
 |\Rm_{k_i(t)}|\le C(\epsilon_i+\sqrt t)/t,\qquad
 \operatorname{inj}_{k_i(t)}\ge c\sqrt t,\qquad
 \|\Rm_{k_i(t)}\|_{2;B_{k_i(t)}(q_i,1/2)}
       \le C(\epsilon_i+\sqrt t).
\]
By its bi-Hölder distance comparison,
a sufficiently small fixed initial ball
\(D_i=B_{k_i(0)}(q_i,r_*)\) is contained in that
time-\(t\) ball for all sufficiently small times.
The other distance inequality puts the image of a
time-\(t\) exponential ball of radius \(c_0\sqrt t\)
inside \(D_i\), with \(c_0>0\) below the uniform
rescaled injectivity radius. These inclusions follow from
the uniformly positive exponents of the bi-Hölder bounds
on protected balls. In particular
\[
 \int_0^t\int_{D_i}|\Ric_{k_i(s)}|^2\,dV_{k_i(s)}ds
       \le C(\epsilon_i^2 t+t^2).
\]
At a point put \(F_t=\int_0^t |\Ric_{k_i(s)}|,ds\). On
the set \(F_t>a\), stop at the first time when the
accumulated integral equals \(a\). The evolving volume
density and the initial density differ by at most a
factor \(e^{2a}\) on that interval because
\(|\partial_s\log dV|=|R|\le2|\Ric|\). By
Cauchy--Schwarz, \(a^2\le t\int|\Ric|^2ds\) on
the stopped interval. Integrate over the exceptional
points and use the spacetime bound. The same argument
from time \(t\) backwards controls the exceptional set
in the time-\(t\) volume. Thus each volume is at most
\(C e^{2a} a^{-2}t^2(\epsilon_i^2+t)\).

Off this set, integrate the evolution equation in
each fixed tangent direction to get
\[
 e^{-2a}k_i(t)\le k_i(0)\le e^{2a}k_i(t).
\]
For the time-\(t\) metric, the sectional curvature after
rescaling is at most \(C(\epsilon_i+\sqrt t)\).
The Jacobi equation in normal coordinates \(\sqrt t\,z\)
gives the same order of coefficient error on the model
ball. The metric comparison gives the initial coefficient
bound outside the exceptional set. Dividing each
exceptional volume by \(t^2\) gives the measure bounds.
\end{proof}

The local result permits exceptional sets with unknown
shape or capacity, and gives no initial first-derivative
control on the fixed model ball. The term
\(\epsilon_i/t\) cannot give a uniform time integral
from zero. Under global upper ratios, diameter, Sobolev
and curvature bounds, Lee and Lee obtain a
Gromov--Hausdorff limit that is a topological manifold
off finitely many points. The stronger small total
curvature case gives a smooth structure via the flow.
First-order convergence of the initial metrics is a
separate requirement
\cite[Theorem 1.3 and Remarks 4.5--4.6]{LeeLee2026Local}.
The harmonic-coordinate route of Deane Yang requires
supercritical local Ricci control
\cite[Theorems 3.1 and 7.1]{Yang1992ConvergenceII}.
For the criterion of Section~\ref{sec:anchored-shell} the
remaining requirements include one wide embedded chart in
the physical source, its center-distance bound, and the
central critical coordinate estimate across any exceptional
regions. Uniform upper physical volume ratios on those scales
are themselves additional to the bounds in
Section~\ref{sec:quantiles}.

% Edward B. Baker III, with OpenAI Codex assistance, 25 November 2026.
% Inputs: repository HEAD 5ff7014 and the dated neck-geometry derivations.
% Exact serving model and reasoning effort were unavailable.

\section{Positive Schouten representatives and boundary charge}
\label{alt:gauge}

Let \(M\) be a smooth closed oriented homotopy four-sphere. Write
\(E=E_+=E_-\) for its common closed chiral Weyl energy, and suppose
\(Y=Y([g])>0\) and \(E\le 8\pi^2-\delta\), \(\delta>0\). For a
representative \(h\) use the Schouten tensor \(P_h=(\operatorname{Ric}_h-R_hh/6)/2\),
and put \(p_h=\sigma_2(h^{-1}P_h)\) and \(T_h=\operatorname{Ric}_{0,h}\).
The elementary symmetric function and Chern--Gauss--Bonnet give
\begin{equation}\label{alt:schouten}
 p_h=\frac{R_h^2}{96}-\frac{|T_h|^2}{8},\qquad
 q:=\int_M p_h\,dV_h=\frac{8\pi^2-E}{2}\ge \frac{\delta}{2}.
\end{equation}
In the unit-volume Yamabe representative \(g_Y\),
\(R_{g_Y}=Y\) and \(Y^2=96q+12\int_M|T_{g_Y}|^2\,dV_{g_Y}\).
In particular \(\sqrt{48\delta}\le Y\le Y_{\mathrm{rd}}=8\sqrt6\pi\).

Theorem A and Corollary B of Chang--Gursky--Yang give a smooth
representative \(h\) with \(R_h>0\) and \(p_h>0\)
\cite[Theorem A, Corollary B and Lemma 1.1]{ChangGurskyYang2002PositiveRicci}.
Their Schouten tensor is \(2P_h\). In the same class one may choose
the unit-volume representative satisfying \(p_h=q\), by their
constant-\(\sigma_2\) existence theorem
\cite[Theorem 1.2 and equation (1.6)]{ChangGurskyYang2003}.
Positivity of \(p_h\) prevents \(R_h\) from vanishing by
\eqref{alt:schouten}. Connectedness and the positive Yamabe
constant, tested against the constant function, give \(R_h>0\).
Uniqueness after volume normalization applies in a fixed nonround
class, and the solutions in the round class are round
\cite[Theorem 1.5]{GurskyStreets2018TwoYamabe}.

For \(S=R_h/6>0\) and least eigenvalue \(a\) of \(h^{-1}P_h\),
Cauchy--Schwarz on the other three eigenvalues gives
\[
 p_h\le a(S-a)+(S-a)^2/3=(S-a)(S+2a)/3.
\]
Thus \(m=S+2a=\lambda_{\min}(\operatorname{Ric}_h)>0\), and
\(p_h\le m(3S-m)/6\le mS/2\). Also \(|P_h|^2<S^2\), because
\(S^2-|P_h|^2=2p_h\). These inequalities give
\begin{equation}\label{alt:riccone}
 R_h\ge\sqrt{96p_h},\qquad
 \lambda_{\min}(\operatorname{Ric}_h)\ge\frac{12p_h}{R_h},
 \qquad \operatorname{Ric}_h<\frac{R_h}{2}h.
\end{equation}
At fixed \(p_h=q>0\) the algebraic bound can degenerate.
The Ricci eigenvalues \(m,(R-m)/3,(R-m)/3,(R-m)/3\), with
\(m=(R-\sqrt{R^2-96q})/4\), have \(m\sim12q/R\). The modified
prescription
\(p_h-\alpha |W_h|_{\mathrm{tens}}^2/16=\lambda\) has
\(\lambda=q-\alpha E/2\ge\delta/4\) at volume one for
\(\alpha=\delta/(16\pi^2)\)
\cite[Theorem 1.1 and equation (1.4)]{ChangGurskyYang2003}.
The same algebraic degeneration is possible with a zero Weyl block.

\subsection{Physical metric estimates}

The chiral Weyl measures and \(|B|\,dV\) are conformally invariant
on each fixed domain. Physical balls and their energy quantile radii
depend on the representative. In any representative one has
\begin{equation}\label{alt:full-curvature}
 \begin{split}
 \int_M |T_h|^2\,dV_h&=\tfrac1{12}\int_M R_h^2\,dV_h-8q,\\
 \int_M |\operatorname{Ric}_h|^2\,dV_h
   &=\tfrac13\int_M R_h^2\,dV_h-8q,\\
 \int_M |\operatorname{Rm}_h|_{\mathrm{tens}}^2\,dV_h
   &=8E-16q+\tfrac13\int_M R_h^2\,dV_h .
 \end{split}
\end{equation}
The first identity follows from \eqref{alt:schouten}. The second uses
the trace decomposition, and the third uses
\(|\operatorname{Rm}|_{\mathrm{tens}}^2
=|W|_{\mathrm{tens}}^2+2|T|^2+R^2/6\).
In the Yamabe gauge \(\int R^2=Y^2\le384\pi^2\). In a positive-Ricci
representative with volume \(V\) and diameter \(D\), Bishop--Gromov gives
\[
 \frac{V}{D^4}r^4\le V_x(r):=\operatorname{Vol}_h B_h(x,r)
 \le \omega_4r^4\quad(0<r\le D),\qquad \omega_4=\pi^2/2.
\]
A scalar upper bound \(R_h\le R_*\) in the constant-\(p_h\) gauge
would give \(D\le\pi\sqrt{R_*/(4q)}\) by \eqref{alt:riccone} and
Bonnet--Myers.

\begin{proposition}[Scalar energy and physical ball bounds]
\label{alt:physical-balls}
Suppose \(h\) has unit volume, \(\operatorname{Ric}_h\ge0\),
\(Y([h])\ge y_0>0\) and \(\int_M R_h^2\,dV_h\le L_*^2\).
Then its diameter is bounded above in terms of \(y_0,L_*\), and
\(c(y_0,L_*)r^4\le V_x(r)\le\omega_4r^4\) for \(0<r\le D\).
For a complete connected noncompact four-manifold with
nonnegative Ricci curvature, finite scalar \(L^2\) energy
and the same Yamabe inequality on compactly supported tests,
the asymptotic volume ratio is at least \((y_0/Y_{\mathrm{rd}})^2\).
\end{proposition}

\begin{proof}
Use a cutoff supported on \(B(x,r)\), equal to one on \(B(x,r/2)\),
with gradient bound \(2/r\). Bishop doubling gives
\(V_x(r/2)\ge V_x(r)/16\). The Yamabe inequality and Hölder yield
\[
 \frac{y_0}{4}\le \frac{24}{r^2}V_x(r)^{1/2}
 +\left(\int_{B(x,r)}R_h^2\,dV_h\right)^{1/2}.
\]
Thus each ball has volume at least \(\kappa r^4\),
\(\kappa=(y_0/192)^2\), or scalar \(L^2\) norm greater than
\(y_0/8\). Pack disjoint unit balls centered two units apart along a
minimizing diameter segment. The two kinds of balls number at most
\(\kappa^{-1}+64L_*^2/y_0^2\) altogether, which bounds \(D\) and gives
the first assertion by Bishop--Gromov.

In the noncompact case let \(\alpha_\infty\) be the asymptotic
volume ratio. Balls of radius \(r\) centered at distance \(2r\) along
a ray have scalar energy tending to zero, which first gives
\(\alpha_\infty\ge\kappa/(81\omega_4)>0\).
Apply the compact-support Yamabe inequality to \(f(d(o,\cdot)/r)\)
for \(f\in C_c^\infty((0,\infty))\). The pushforward of
\(r^{-4}dV\) by \(d(o,\cdot)/r\) tends on compact subintervals to
\(\alpha_\infty\,d(\omega_4s^4)\), by convergence of the radial volume
distribution. The distance gradient has norm one almost everywhere.
The scalar term divided by \(r^2\) tends to zero by Hölder and
finiteness of the scalar energy on the escaping supports. The limit
inequality bounds \(y_0\) by \(\sqrt{\alpha_\infty}\) times the
Euclidean Yamabe quotient of the radial test. Cutoffs of the
stereographic round conformal factor have quotients tending to
\(Y_{\mathrm{rd}}\) \cite[\S3, Theorems 3.2--3.3]{LeeParker1987Yamabe}.
\end{proof}

For a supplied complete noncompact smooth pointed limit at
vanishing scales in these constant-\(p_h\) gauges, assume the
uniform scalar-energy bound of Proposition~\ref{alt:physical-balls}.
Nonnegative Ricci curvature and the compact-support Yamabe
inequality pass to the limit. The
scalar \(L^2\) budget passes on exhaustion and the constant \(p_h\)
rescales to zero. Its volume ratio is at least
\(\delta/(8\pi^2)\). A single smooth ALE end with link \(S^3\)
has volume ratio one, in which case Bishop--Gromov rigidity gives
Euclidean space.

Positivity alone leaves the scalar-energy and diameter estimates
open. For example, the rotational ellipsoids
\(z^2/L^2+|y|^2=1\) in \(\mathbb R^5\) are strictly convex and
conformally round. In the warped-product expression
\(ds^2+a(s)^2\gamma_{S^3}\), this follows by using
\(d\log\rho=ds/a(s)\), with smooth pole expansions
\(a(s)=s+O(s^3)\) in pole distance. If \(k_r,k_t\) are radial
and tangential sectional curvatures, then
\(p_h=3k_rk_t/2>0\). On \(|z|\le L/2\),
\(a=\sqrt{1-z^2/L^2}\) and \(k_t=(1-a_s^2)/a^2\) have uniform
positive lower bounds. The volume is comparable to \(L\) and
\(\int R^2\ge cL\). Unit-volume rescaling leaves that divergence
and gives diameter at least \(2L/V^{1/4}\). A unit ball in that
rescaling lies in an original \(z\)-slab of width \(2V^{1/4}\),
so has unit-volume measure at most \(CV^{-3/4}\). This family has
nonconstant \(p_h\). The constant-\(p_h\) choice in its round class
is round. For varying constant-\(p_h\) classes,
\eqref{alt:full-curvature} and Proposition~\ref{alt:physical-balls}
use a further uniform scalar-energy assumption. Transferring a
physical radius selection between gauges also requires distance
comparison or a new selection in the chosen metric.

\subsection{Conformal boundary current and its modular consequence}

Fix an actual compact smooth domain \(U\subset M\).
Let \(h=e^{2u}g\), \(g_t=e^{2tu}g\), and
\(N_t=(\operatorname{tr}_{g_t}P_{g_t})g_t-P_{g_t}\).
The contracted Bianchi identity gives \(\operatorname{div}N_t=0\).
The infinitesimal Schouten law
\(\partial_tP_{g_t}=-\nabla^2_{g_t}u\) and
\(\operatorname{tr}(T_1(K)K)=2\sigma_2(K)\) for
\(K=g_t^{-1}P_{g_t}\) give
\begin{equation}\label{alt:newton}
 \partial_t(p_{g_t}\,dV_{g_t})
 =-\operatorname{div}_{g_t}(N_t^{\sharp_{g_t}}\nabla^{g_t}u)\,dV_{g_t}.
\end{equation}
Here the term \(-4up_{g_t}\) from differentiating the inverse
metric cancels the volume derivative. These are the Newton tensor
identities of \cite[\S2.1, equations (2-1)--(2-5) and Lemma 2.1]{GurskyStreets2018TwoYamabe}
with the sign appropriate to \(g_t=e^{2tu}g\). Write
\[
 q_g(U)=\int_U p_g\,dV_g,\qquad
 J=\sum_{\Sigma\subset\partial U}
   \int_0^1\!\int_\Sigma N_t(\nu_t,\nabla^{g_t}u)\,dS_{g_t}\,dt,
 \qquad q_h(U)-q_g(U)=-J ,
\]
using the outward \(g_t\)-unit normals. In this integral the changes
of the normal, gradient and area contribute factors
\(e^{-tu},e^{-2tu},e^{3tu}\), whose product is one. Integrating
\(P_{g_t}=P_g-t\nabla_g^2 u+t^2du\otimes du
-t^2|du|_g^2g/2\) gives the exact face current
\begin{equation}\label{alt:face-current}
 J_\Sigma=\int_\Sigma\!
  \left[N_0(\nu,\nabla u)
   +\tfrac12\bigl(\nabla^2u(\nu,\nabla u)-(\Delta_g u)u_\nu\bigr)
   -\tfrac12|\nabla u|^2u_\nu\right]dS_g .
\end{equation}
Here \(\Delta_g=\nabla^i\nabla_i\). If \(f=u|_\Sigma\),
\(v=\nabla_\Sigma f\), \(a=u_\nu\), and
\(II(X,Y)=g(\nabla_X\nu,Y)\), the same integral has integrand
\[
 aN_0(\nu,\nu)+N_0(\nu,v)-a\Delta_\Sigma f
 -\tfrac12 II(v,v)-\tfrac12(\operatorname{tr}II)a^2
 -\tfrac12a(a^2+|v|^2).
\]
Indeed the normal-normal Hessian cancels and the mixed Hessian
contributes \(\langle\nabla_\Sigma a,v\rangle-II(v,v)\).
The two integrals of the tangential Laplacian combine by integration
by parts. Hölder bounds \(|J_\Sigma|\) by
\[
 C\bigl(A_\Sigma D_\Sigma+\|II\|_3 A_\Sigma^2+A_\Sigma^3\bigr),
 \quad A_\Sigma=\|\,|a|+|v|\,\|_3,\quad
 D_\Sigma=\|N_0(\nu,\cdot)\|_{3/2}+\|\Delta_\Sigma f\|_{3/2}.
\]
All the norms in this bound are invariant under constant rescaling.
For paired faces with a common forward transverse direction,
\eqref{alt:face-current} gives the difference of the two integrals
in that direction. Cancellation of the signed sum is possible even
when individual integrals are large.

For either chirality put
\(I_\pm(U,h)=E_\pm(h;U)+2q_h(U)\), and denote by
\(C_\pm(\partial U,h)\) the oriented sum of the Chern--Simons
integrals of the physical Levi-Civita half-spin connections in
boundary frames. Use the unnormalized anti-Hermitian fundamental
trace and set \(P_0=8\pi^2\), \(s_+=-1\), \(s_-=1\).
The homotopy sphere is spin because \(H^2(M;\mathbb Z/2)=0\).
Its half-spin bundles restrict trivially to the boundary
since \(\pi_j(\mathrm{SU}(2))=0\) for \(j=0,1,2\).
The diagonal Riemann curvature block is
\(W^\pm+R_h\operatorname{Id}/12\) and the off-diagonal block
has squared norm \(|T_h|^2/4\). The factor \(1/2\) in the
Clifford curvature map is paired with
\(-\operatorname{tr}(c(\omega)c(\eta))=4\langle\omega,\eta\rangle\)
for two-forms of the acting chirality
\cite[\S2.3]{Baker2026DiracYamabe}. Evaluating
\(s_\pm\operatorname{tr}(F_\pm\wedge F_\pm)\) gives the difference
of their squared norms times \(dV_h\), namely
\((|W^\pm|^2+R_h^2/48-|T_h|^2/4)dV_h\).
Subtracting the boundary transgression gives the relative second
Chern integer, so
\[
 I_\pm(U,h)=s_\pm\bigl(P_0k_\pm+C_\pm(\partial U,h)\bigr),
 \qquad k_\pm\in\mathbb Z .
\]
The transgression identity used here is
\cite[\S3, Proposition 3.2]{ChernSimons1974Characteristic}.
Using the canonical conformal spin isometry to transport the frames
preserves \(k_\pm\). Hence
\begin{equation}\label{alt:gauge-charge}
 I_\pm(U,h)-I_\pm(U,g)=-2J,\qquad
 C_\pm(\partial U,h)-C_\pm(\partial U,g)=-2s_\pm J .
\end{equation}
Independent frames give the second equation modulo \(P_0\).
For a single face the real equality follows by applying the
formula to a collar with the conformal change cut off at its
other end. In particular, using \(d_P(x)\) for the distance
to \(P\mathbb Z\), one gets
\[
 d_{P_0}\bigl(C_\pm(\partial U,h)\bigr)
 \le d_{P_0}\bigl(C_\pm(\partial U,g)\bigr)+2d_{P_0/2}(J).
\]
For example, if \(p_h\ge0\) on \(U\), then
\(0\le E_\pm(h;U)\le |I_\pm(U,g)|+2\max\{-J,0\}\).
In a volume-one Yamabe source the positive part of \(J\) is
bounded by \(Y^2\operatorname{Vol}_g(U)/96\) by \eqref{alt:schouten}.
Controlling the other sign is an additional boundary estimate.

Global curvature and volume bounds alone leave that sign open
on an actual residual. On the unit round sphere let
\(h_\lambda=\Phi_\lambda^*g\), where in stereographic coordinates
\(\Phi_\lambda(x)=\lambda x\), and take
\(U_\lambda=\{a/\lambda\le |x|\le b/\lambda\}\) for fixed
\(0<a<b\). The two gauges have the same global curvature bounds
and positive constant \(p=3/2\). The physical \(g\)-volume of
\(U_\lambda\) tends to zero and its \(h_\lambda\)-volume is the
fixed round volume of \(\{a\le |x|\le b\}\). Thus
\(J=q_g(U_\lambda)-q_{h_\lambda}(U_\lambda)\) has a negative
nonzero limit. This example concerns the zero Weyl class.

\begin{proposition}[Positive chiral charge on a residual domain]
\label{alt:positive-charge}
If \(p_h>0\) globally, put \(U^c=\overline{M\setminus U}\).
Then
\[
 d_{P_0}\bigl(C_\pm(\partial U,h)\bigr)
 =\min\{I_\pm(U,h),I_\pm(U^c,h)\}
 \ge\min\{E_\pm(h;U),E_\pm(h;U^c)\}.
\]
In particular, along a sequence with complementary chiral mass
bounded below and modular distance tending to zero, the chiral
Weyl energy and the Schouten integral on \(U\) both tend to zero.
\end{proposition}

\begin{proof}
The two charges are nonnegative and sum to \(P_0\) by
\eqref{alt:schouten}. The transgression congruence gives the
minimum exactly. If \(I_\pm(U^c,h)\) is bounded below by a positive
constant, the smaller charge is eventually \(I_\pm(U,h)\).
In the constant-\(p_h\) gauge,
\[
 I_\pm(U,h)=E_\pm(h;U)+(P_0-E)
                \frac{\operatorname{Vol}_h(U)}{\operatorname{Vol}_h(M)}.
\]
Thus the volume fraction is bounded by the modular distance divided
by \(\delta\) once the smaller charge is selected.
The same argument works in an iterated exhaustion
\(\lim_{L\to\infty}\limsup_i\), provided the complementary lower
bound holds in \(\liminf_i\) with one constant for all sufficiently
large fixed \(L\).
\end{proof}

A single balanced cut has modular distance at least the smaller of
its two Weyl masses. For a middle residual one needs the oriented
sum on its actual boundary. On a radial spherical collar
\(h=e^{2w}(dt^2+\gamma_{S^3(1)})\), let \(v=w'\).
The radial and tangential sectional curvatures are
\(-e^{-2w}v'\) and \(e^{-2w}(1-v^2)\). Since \(W=0\) and
\(p_h=3k_rk_t/2\), integration gives
\[
 I_\pm([t_-,t_+]\times S^3,h)=F(v(t_+))-F(v(t_-)),\qquad
 F(v)=4\pi^2-6\pi^2v+2\pi^2v^3.
\]
For the slice with increasing-\(t\) coorientation the modular
Chern--Simons value is \(s_\pm F(v)\). To fix the additive
constant, fill the slice from below by a radial ball whose metric
agrees near the slice and has \(w=t+\mathrm{const}\) at the
early pole. That pole has slope one and \(F(1)=0\).
For \(|v|\le1\) the modular distance on the slice is
\(2\pi^2(1-|v|)^2(2+|v|)\). Each boundary of an exact round
product cylinder has distance \(4\pi^2\), while the oriented
sum is zero. Strictly admissible round-class examples with
the same limiting cancellation are supplied by
\[
 w_L(t)=L-\log 2-\tfrac12\bigl(\log\cosh(t+L)+\log\cosh(t-L)\bigr).
\]
Here \(-1<v_L<1\), \(v_L'<0\), both pole expansions are smooth,
and \(w_L,v_L\to0\) uniformly on \([-L/2,L/2]\). The round
strip \([-L,0]\) with \(w=-\log\cosh t\) instead has limiting
charge \(4\pi^2\). The symmetric strip \([-L,L]\) in that
metric has charge tending to \(P_0\), so its modular distance
vanishes on the complementary branch. These examples explain
the paired orientation and the uniform complementary mass
hypothesis of Proposition~\ref{alt:positive-charge}.
For example, in a fixed admissible perturbation of the round
metric with nonzero Weyl tensor, take \(U\) to be the exterior
of a ball shrinking to a point with nonzero curvature in the
tested block. The complementary mass is positive and tends to
zero, while the exterior Weyl mass tends to the positive mass
of the whole sphere.

Even the local equation \(p_h=q>0\) permits two large individual
fluxes with a small oriented sum. On \([-L,L]\) set
\(w(0)=-\log L\), \(v(0)=0\) in
\[
 w'=v,\qquad v'=-\tfrac{2q}{3}e^{4w}/(1-v^2).
\]
As long as \(|w+\log L|\le1\) and \(|v|\le1/2\),
\(|v'|\le C_q L^{-4}\). Integration closes these bounds for
large \(L\) and gives \(|v|\le C_qL^{-3}\) and
\(|w+\log L|\le C_qL^{-2}\) on the full interval.
The individual modular distances tend to \(4\pi^2\), and their
paired charge is \(O_q(L^{-3})\). In this solution the first
integral \((1-v^2)^2-\tfrac{2q}{3}e^{4w}\) tends to one as
\(L\to\infty\). A smooth pole in a global radial completion
with constant \(q\) would force that first integral to be zero.

\section{Round-base restart and the Green-rescaling budget}
\label{alt:base}

Let \(I^+(M)\) be the infimum of \(E_+\) over smooth
positive-Yamabe classes on a fixed oriented smooth homotopy
four-sphere \(M\). Consider volume-one Yamabe minimizers \(g_i\)
with common chiral energy \(E_i\in[4\pi^2,8\pi^2-\delta]\),
\(E_i\to e\), \(Y_i\to y>0\), and
\(\mathfrak B(g_i):=\int_M|B_{g_i}|\,dV_{g_i}\to0\).
We first use a physical comparison with a round base, then a
comparison in a potentially different conformal gauge.

\subsection{A restart at one physical puncture}

\begin{theorem}[Restart on the original smooth type]\label{alt:restart}
Suppose \(M\) is nonstandard, \(e=I^+(M)\) and
\(y=Y_{\mathrm{rd}}\). Supply a volume-one round sphere
\((N,q_N)\) with a finite marked set \(P\), and oriented embeddings
on neighborhoods of whole fixed compact truncations of
\(N\setminus P\) on which the physical metrics converge smoothly
to \(q_N\). Then there is a slow coherent exhaustion, and after
a subsequence one fixed puncture label, whose actual exterior
contains \(e-o(1)\) of both chiral Weyl masses. Its intrinsic
physical diameter tends to zero. Capping that exterior and
selecting a volume-one Yamabe minimizer gives metrics \(g_i'\)
on the original oriented smooth type with
\[
 E_+(g_i')\to I^+(M),\qquad
 \mathfrak B(g_i')\to0,\qquad Y([g_i'])\to Y_{\mathrm{rd}} .
\]
Each restarted class has a marked conformally round germ.
\end{theorem}

\begin{proof}
A comparison on the whole unpunctured sphere would imply
\(E_i\to0\), so \(P\ne\varnothing\). Use disjoint round normal
coordinate balls with radius variable \(\rho\) around its points.
At cut radius \(r\), set \(t=\log(r/\rho)\) and choose the angular
orientation consistently. In the \(\rho^{-2}\) gauge the round
comparison is
\[
 dt^2+\left(\frac{\sin(k\rho)}{k\rho}\right)^2\gamma_{S^3(1)},
\]
where \(k^2>0\) is the round sectional curvature. The product
\(C^2\) error through \(|t|\le\ell+2\) is
\(O((r e^{\ell+2})^2)\). The \(C^4\) error on the central unit
splices is \(O(r^2)\). At fixed \(r,\ell\) the smooth physical
comparison gives an additional \(o_i(1)\) in both estimates.

Choose lengths \(\ell_m\to\infty\), tolerances
\(\eta_m\downarrow0\) with \(\ell_m\eta_m^2\to0\), then radii
\(r_m\downarrow0\) small enough for the buffered closures to be
disjoint and the round errors below \(\eta_m/2\). Also require
\(r_m^4(1+\ell_m)\to0\). Choose increasing source thresholds at which
the errors of the source on the entire larger truncation
\(N\setminus\bigcup_{p\in P} B_{q_N}(p,r_m e^{-\ell_m-2})\)
are below \(\eta_m\) and the volume integral over the truncation
at \(r_m\) agrees with its round value to within \(\eta_m\).
Impose unscaled smooth comparison through order \(m+4\) at each
threshold with error tending to zero. At each index use only the
map for the latest threshold and its restrictions. Write
\(r_i',\ell_i'\) for the resulting cut radius and buffer length.
This choice supplies the nested truncations and their comparisons
on all fixed compact subsets.

By the whole-truncation topology in the proof of
Theorem~\ref{thm:finite-punctures}, there is one exterior \(V_{ij}\)
per face and its marked spherical cap is a homotopy sphere.
Capping the central truncation gives \(N\), and the marked
connected sum of all capped pieces gives \(M\).
Let \(m^\pm_{ij}=E_\pm(g_i;V_{ij})\) and let \(e^\pm_{ij}\)
be the capped chiral energies. The buffered cap estimates in
Theorem~\ref{thm:caps} give
\begin{equation}\label{alt:restart-cap}
 \begin{split}
 |e^\pm_{ij}-m^\pm_{ij}|&\le C\eta_{2,i}^2,\\
 \mathfrak B(\widehat g_{ij})&\le
         \mathfrak B(g_i;V_{ij})+C\eta_{4,i},\\
 Y([\widehat g_{ij}])&\ge
   \min\{Y_i,(1-C\eta_{2,i})Y_{\mathrm{rd}}\}
        -C(\ell_i')^{-3/2}.
 \end{split}
\end{equation}
Here \(\eta_{2,i},\eta_{4,i}\le\eta_m\) along the chosen
exhaustion. The first bound compares both old and new unit splices
to the conformally flat product. The second uses the fourth jets
on the central splice and the conformal invariance of the Bach
\(L^1\) density in dimension four. Both bounds are independent of
\(r_i'\). One may put the splice and cap at physical length
scale \(r_i'\), using the factor \(r_i'/\rho\) on the retained
collar, extended with a cutoff on the spare buffer. The cap in this
pre-minimization representative has volume \(O((r_i')^4)\)
and the replaced representative on the long collar has volume
\(O((r_i')^4(1+\ell_i'))\).

All capped exteriors have positive Yamabe constant with limits
\(Y_{\mathrm{rd}}\), by \eqref{alt:restart-cap} and the universal
upper bound. The central capped piece carries the round metric.
By Theorem~\ref{thm:nonsplitting}, the strict upper band and the
\(o(1)\) cap costs imply at most one capped piece of energy
\(\ge4\pi^2\). The nonstandard type forces one nonstandard
capped exterior, of the original oriented type. For its label
\(j_i\),
\[
 I^+(M)\le e^+_{i,j_i},\qquad
 E_i-m^+_{i,j_i}\le E_i-I^+(M)+C\eta_{2,i}^2\to0 .
\]
The minus estimate follows because both closed energies are equal
on each capped piece and on each source. Fix the surviving label
along a subsequence. Larger cuts in the same normal coordinate
range, defined by restriction of the same map, add a marked
cylindrical layer. The capped type is preserved and the larger
exterior still contains the limiting mass. Smaller cuts within
the controlled buffer lose at most
\(C(1+\ell_i')\eta_{2,i}^2=o(1)\) of either Weyl mass.

The exteriors have physical volume tending to zero by the
volume-one normalization and the imposed convergence of the
central volume. The source Yamabe inequality gives uniform
physical lower bounds \(\operatorname{Vol} B(x,u)\ge c_0u^4\)
for \(0<u\le r_0\). In fact the cutoff in a ball and the
bounded source scalar curvature give \(V_x(u)/u^4\ge
c(V_x(u/2)/(u/2)^4)^{1/2}\). Iteration to small radii gives
the uniform bound. For each exterior point, a ball of radius
\(\min\{d(x,\partial V_{ij})/2,r_0\}\) lies in its exterior.
Its small volume thus implies
\(d(x,\partial V_{ij})\le 2c_0^{-1/4}
\operatorname{Vol}(V_{ij})^{1/4}\) for large \(i\).
Shortest segments to the boundary stay in the exterior up to the
first boundary point. Its sole boundary sphere has intrinsic
diameter \(O(r_i')\), giving the intrinsic diameter bound
\[
 \operatorname{diam}_{\rm intr} V_{ij}
 \le 4c_0^{-1/4}\operatorname{Vol}(V_{ij})^{1/4}+Cr_i'\to0 .
\]

Transport the surviving capped class to \(M\) by the oriented
diffeomorphism and take a unit-volume Yamabe minimizer
\cite[\S3]{Schoen1984Yamabe}. The class invariants are unchanged.
Equation \eqref{alt:restart-cap} and the minimizing bound imply the
three asserted limits, and the exact cap supplies the marked germ.
\end{proof}

For a given \(i\) the signed change equals the splice excess
\(e^+_{i,j_i}-m^+_{i,j_i}\) minus the discarded mass
\(E_i-m^+_{i,j_i}\). A strict improvement requires the second
term to exceed the first. A further restart requires new usable
comparison data. An estimate for the size of the marked germ
in the restarted Yamabe gauge requires further control.

\subsection{A conformal Green model}

For the same source hypotheses preceding
Theorem~\ref{alt:restart}, allow general \(y>0\) and \(e\).
Supply a smooth closed oriented homotopy-sphere conformal base
\((N,h)\) with finite marked set \(P\). Suppose on full compact
truncations off \(P\) there are oriented comparison embeddings
\(F_i\) and representatives \(F_i^*g_i=u_i^2h_i\), where
\(h_i\to h\) smoothly on fixed compacts. Compact-support Yamabe
tests then give \(Y([h])\ge y\) by cutoff at the finitely many
smooth points. The invariance of the Bach \(L^1\) density gives
\(B_h=0\) away from \(P\), and hence everywhere by smoothness.
Write \(e_N=E_\pm(h)\le e\), by conformal measure invariance
and exhaustion.

\begin{theorem}[One-pole Green budget]\label{alt:green}
Suppose a further constant length normalization at scales
\(\lambda_i\downarrow0\) gives
\[
 \lambda_i^{-2}F_i^*g_i\longrightarrow g_X=w^2h,\qquad
 w=\sum_{j=1}^{\ell}a_jG_{p_j},\quad a_j>0,
\]
smoothly on fixed compact subsets of \(N\setminus P\).
Here \(p_j\in P\) are distinct poles and \(G_p\) is the Green
function satisfying
\((6\nabla^*\nabla+R_h)G_p=24\pi^2\delta_p\).
On \(X=N\setminus\{p_1,\ldots,p_\ell\}\), including any ordinary
marked points, one has
\[
 \int_X |\operatorname{Ric}_{g_X}|^2\,dV_{g_X}
       =4e_N+32\pi^2(\ell-1),\qquad
 e-e_N\ge 8\pi^2\ell-y^2/48 .
\]
There is exactly one pole in the strict upper band, and hence
\(e-e_N\ge(Y_{\mathrm{rd}}^2-y^2)/48\).
\end{theorem}

\begin{proof}
A positive homogeneous solution with no pole is impossible on the
closed base with \(Y([h])>0\). The conformal scalar-curvature law
gives \(R_{g_X}=0\) away from the poles. Near each pole,
the corresponding Green function has positive \(r^{-2}\) leading
term in conformal normal coordinates
\cite[Lemma 6.4 and Theorem 6.5]{LeeParker1987Yamabe}. The other
terms are smooth. Inversion and a constant coordinate rescaling
give one Euclidean, unquotiented end per pole with differentiable
decay \(O_2(R^{-\alpha})\) for any fixed \(0<\alpha<2\).
The metric is complete and its curvature square is integrable
by the decay. A large end truncation removes \(\ell\) balls from
\(N\), so has Euler characteristic \(2-\ell\).
The Gauss--Bonnet boundary integral tends to \(8\pi^2\) on each
end: the rescaled end collar tends in \(C^2\) to the Euclidean
unit sphere collar with outward normal. The sign and coefficient
are those of a Euclidean ball
\cite[\S7, equation (7.2)]{TianViaclovsky2005ALE}. Since the
total chiral Weyl energies on \(X\) equal \(e_N\), the open
Gauss--Bonnet identity is
\[
 8\pi^2(2-\ell)=2e_N-\tfrac12
       \int_X|\operatorname{Ric}_{g_X}|^2\,dV_{g_X}
       +8\pi^2\ell .
\]

On an embedded fixed compact away from all of \(P\), the limit
Ricci integral equals the limit of physical trace-free Ricci
integrals, using only constant-scale invariance and smooth
convergence to a scalar-flat metric. Exhaustion proves the bound
by the limiting source budget
\[
 \lim_i\int_M|\operatorname{Ric}_{0,g_i}|^2,dV_{g_i}
     =4e+y^2/12-32\pi^2
\]
from \eqref{alt:schouten}. The ordinary omitted marks carry
no smooth curvature mass. The two Ricci identities give the
asserted inequality. Two poles would force \(e\ge e_N+8\pi^2\)
since \(y\le Y_{\mathrm{rd}}\), contradicting the strict band.
\end{proof}

Here the Green model and the compact conformal base have the same
Weyl measures on the comparison domains. A separate disjoint bubble
would need additional comparisons. If \([h]\) is nonround,
Theorem~\ref{thm:closed-gap} gives \(e_N>4\pi^2\) and
\(y\le Y([h])<Y_{\mathrm{rd}}\) by strict Yamabe rigidity
\cite[Theorem 2]{Schoen1984Yamabe}. Since
\(e-e_N<4\pi^2-\delta\), Theorem~\ref{thm:finite-punctures}
gives the same smooth type as \(M\) from the conformal comparison.
Thus \(I^+(M)\le e_N\), and the Green model is excluded at the
infimum. More generally the uniform exact-class Yamabe deficit
\(Y([h])\le Y_{\mathrm{rd}}-d_\delta\) in this closed band
gives
\[
 e-I^+(M)\ge e-e_N\ge
 \frac{Y_{\mathrm{rd}}^2-(Y_{\mathrm{rd}}-d_\delta)^2}{48}>0 .
\]
This uses the deficit on the closed base in
Lemma~\ref{thm:exact-compactness} and
\cite[\S12.2]{Baker2026TwistorSections}.

For a nonround base, the supplied zero-factor comparisons can be
normalized by the factor at an ordinary point. Harnack bounds and
interior regularity give a positive limiting factor solving the
homogeneous conformal Laplace equation away from \(P\). The
nonnegative isolated point sources give the Green representation
in Theorem~\ref{alt:green} under the smooth coherent comparison
hypotheses \cite{Baker2026MarkedLimits}. A positive unrescaled
limit factor instead extends through the marked points. Its
volume \(v\le1\) satisfies \(Y([h])v^{1/2}\le yv\) by the
Yamabe inequality and the limiting equation, so \(v=1\) and
\(Y([h])=y\). The physical metrics converge in the measured
Gromov--Hausdorff sense along that subsequence by whole-truncation
comparison, collapse of the one-face complements and source
ball noncollapse \cite{Baker2026MarkedLimits}.
In this physical branch at the infimum all chiral,
scalar-square, trace-free Ricci and full curvature defects vanish
by volume and energy saturation and \eqref{alt:full-curvature}.
For a round base, if the Green comparison is supplied then
\(e_N=0\) and the single-pole Ricci identity makes the complete
scalar-flat conformally flat model Euclidean. Other ordinary
marks can still be omitted by the comparison. The resulting
flat model carries zero chiral mass and the source may retain
energy outside the comparison.

\section{Closed critical metrics}
\label{alt:closed}

\subsection{Isolation of the first exact critical value}

For a fixed smooth closed oriented homotopy four-sphere \(M\), let
\[
 \mathcal A_M=\{E_+([g]): B_g=0,\ Y([g])>0,\ 0<E_+([g])<8\pi^2\}
\]
range over smooth metrics on \(M\), and set
\(a_M=\inf\mathcal A_M\) when the set is nonempty. Let
\(\kappa_\delta>0\) be the exact same-type energy drop in the
upper band \(E_+\le 8\pi^2-\delta\) from
Theorem~\ref{thm:ale-descent}.

\begin{theorem}[First fixed-type critical value]\label{alt:critical}
If \(\mathcal A_M\ne\varnothing\), then its first value is attained
and \(4\pi^2<a_M<8\pi^2\). For
\[
 a_M\le b\le8\pi^2-\delta,\qquad b<a_M+\kappa_\delta ,
\]
the volume-one Yamabe representatives realizing values in
\(\mathcal A_M\cap[a_M,b]\) are sequentially compact up to oriented
diffeomorphisms. The set of distinct values in that sublevel is
finite. In particular \(a_M\) is isolated among exact positive
subcritical values on \(M\).
\end{theorem}

\begin{proof}
Take a sequence of exact values converging to \(e\le b\).
Theorem~\ref{thm:closed-gap} places the representatives in the
exact compactness band. In the blow-up alternative of
Lemma~\ref{thm:exact-compactness}, the smooth ALE compactification
and the marked same-type descent produce a member of
\(\mathcal A_M\) of value at most
\(e-\kappa_\delta<a_M\), which is impossible. The compact
alternative gives smooth convergence with limiting energy \(e\)
and positive Yamabe sign. Apply this first to a minimizing
sequence, taking \(\delta>0\) with \(a_M<8\pi^2-\delta\) and
discarding a finite prefix, to attain \(a_M\). The closed gap
gives the strict lower bound at the attaining class.

For finiteness, fix any smooth Bach-flat metric \(g\) on \(M\)
and take the trace-free Bach gauge of Gursky--Viaclovsky
\cite[\S2.3, equation (2.59) and Theorem 2.11]{GurskyViaclovsky2015}.
For \(0<\alpha<1\) its map is
\[
 P_g^B(\theta)=\pi_0^g\!
     \left(B_{g+\theta}
         +K_{g+\theta}(\square_{K_g}\beta_g\theta)\right)
 : C^{4,\alpha}(S^2_0)\longrightarrow C^{0,\alpha}(S^2_0)
\]
with the normalization of that source. Here \(K_g\) is its conformal
Killing operator, \(\square_{K_g}=\delta_g K_g\),
\(\beta_g\theta=\delta_g\theta\) on trace-free tensors, and \(\pi_0^g\)
is the fixed-background trace-free projection. This map is real
analytic in \(\theta\): inverse metric, differentiation with fixed
background, and multiplication through the required jet orders
are analytic on the open cone of positive metrics. Theorem 2.11
gives an elliptic derivative and puts every sufficiently
\(C^{4,\alpha}\)-close Bach-flat metric into the small zero set by
a diffeomorphism near the identity and a conformal change.
The elliptic derivative has a finite-dimensional kernel \(K\)
and closed image of finite codimension. Project onto that image
and apply the analytic implicit function theorem on a closed
complement of \(K\). Its graph \(\theta=k+v(k)\) contains all
small zeros.

The full four-index Weyl integral \(F(k)\) on the graph is
real analytic by the same metric operations and the analytic
volume density. At a parameter from an actual Bach-flat metric
its unrestricted first variation vanishes, so \(dF(k)=0\) on
the entire kernel tangent space. The finite-dimensional
\L{}ojasiewicz gradient inequality
\cite[Theorem 1]{Feehan2020Gradient} gives
\(|F(k)-F(0)|^\vartheta\le C|dF(k)|\), with
\(0<\vartheta<1\) near zero. Consequently all nearby smooth
Bach-flat metrics have the same full Weyl integral, and hence
the same \(E_+\) on \(M\). If the displayed sublevel contained
infinitely many distinct values, choose a representative
of each and use smooth subsequential compactness. Local constancy
at its limit contradicts distinctness.
\end{proof}

The isolation concerns values. A least exact critical value
may be a saddle for variations through arbitrary positive-Yamabe
metrics. More generally, distinct exact values accumulating
inside a fixed band force the blow-up alternative along the
corresponding sequence, and an exact value lower by at least
\(\kappa_\delta\) than the accumulation value on the same type.

\subsection{The chiral Cotton obstruction and Einstein reduction}

Let \(C_{ijkl}=(W^+)_{ijkl}\) be the all-lowered Weyl block,
\(D_{jkl}=\nabla^i C_{ijkl}\), and
\(L(\alpha)_{jkl}=\alpha^iC_{ijkl}\) for a real one-form.
Put \(w=|W^+|_{\mathrm{HS}}\) and \(U=\{w>0\}\).
Use the full ordered-index norm for three-tensors.
In an orthonormal self-dual eigenbasis \(\Sigma_a\), with
eigenvalues \(\lambda_a\),
\(C_{ijkl}=\sum_a\lambda_a(\Sigma_a)_{ij}(\Sigma_a)_{kl}\).
Contraction of each normalized eigenform with a vector has
squared norm half the vector's squared norm, and summing
over the ordered last pair contributes a factor of two.
Thus \(L^*L=w^2\operatorname{Id}\). Define on \(U\)
\[
 \beta=-w^{-2}L^*D,\qquad
 \mathcal R=D+L(\beta),\qquad r=|\mathcal R|,
\]
and set \(\mathcal R=D\) on \(M\setminus U\).
The conformal change \(\widehat g=e^{2f}g\) gives
\(\widehat C=e^{2f}C\) and \(\widehat D=D+L(df)\) as covariant
tensors \cite[equations (3.3)--(3.4)]{GurskyMcKeownTyrrell2026}.
Because \(\widehat L=L\), \(\widehat L^*=e^{-4f}L^*\), and
\(\widehat w^2=e^{-4f}w^2\), the residual \(\mathcal R\) is
unchanged and \(\widehat\beta=\beta-df\). The pointwise orthogonal
decomposition is
\begin{equation}\label{alt:cotton-split}
 |D+L(df)|_g^2=r^2+w^2|df-\beta|_g^2 \quad\text{on } U.
\end{equation}
The left-hand side is \(r^2\) on the complement.
In particular
\[
 \int_M |\widehat D|_{\widehat g}^{4/3}\,dV_{\widehat g}
 \ge\int_M r^{4/3}\,dV_g ,
\]
using the three-index norm's conformal weight. On any piecewise
smooth regular closed curve \(\gamma\subset U\), the integral
of the left side of \eqref{alt:cotton-split} with respect to
arclength is bounded below by
\[
 \int_\gamma r^2\,ds_g+
         \frac{|\oint_\gamma\beta|^2}{\int_\gamma w^{-2}\,ds_g}.
\]
This follows from the tangential component of \(df-\beta\),
Cauchy--Schwarz and \(\oint_\gamma df=0\).

For a smooth half-harmonic representative, \(\mathcal R=0\)
on \(U\) and \(\beta\) must extend to a smooth closed one-form.
Conversely such an extension is exact on a homotopy sphere.
The resulting divergence equation extends by continuity to
\(\overline U\), and holds on the interior of the complement
because \(C\) vanishes there.

\begin{theorem}[Closed harmonic scale for a Bach-flat class]
\label{alt:harmonic-scale}
Let \(g\) be smooth and Bach-flat on a closed oriented homotopy
four-sphere and suppose \(W^+\not\equiv0\). Then a smooth
half-harmonic representative for \(W^+\) exists if and only if
\(\mathcal R=0\) on \(U\) and \(\beta\) has a smooth extension
to the whole sphere. Every such half-harmonic representative
is Einstein.
\end{theorem}

\begin{proof}
With the Schouten tensor \(P_g\), the Bach formula in the
\(W^+\) block is
\(B_{ij}=2(\nabla^k\nabla^l+P_g^{kl})C_{ikjl}\)
\cite[\S3, p.~78]{LeBrun2013EinsteinExtremal}.
Use \(\nabla^l C_{ikjl}=-D_{jik}\) and substitute
\(D=\mathcal R-L(\beta)\) twice to obtain
\begin{equation}\label{alt:cotton-bach}
 \tfrac12 B_{ij}
 =C_{ikjl}\bigl(P_g^{kl}-\nabla^k\beta^l+\beta^k\beta^l\bigr)
       -\nabla^k\mathcal R_{jik}-\beta^k\mathcal R_{ikj}.
\end{equation}
Here the derivative of the Weyl tensor in the substitution
uses \(\nabla^k C_{ajik}=-D_{iaj}\).
When \(\mathcal R=0\), the skew part yields
\(W^+((d\beta)^+)=0\) by the Bianchi identity.

The invertible locus of \(W^+\) is dense in \(U\) when
\(D=-L(\beta)\) there. Indeed, if rank two held on an open
subset, a local quaternionic frame \(I,J,K=IJ\) could be
chosen with \(J\) the zero eigenline, and with opposite
eigenvalues on \(I,K\). Differentiate \(W^+(\omega_J)=0\)
and contract on the first index of \(C\). The contracted
derivative of \(C\) vanishes against \(\omega_J\) by
\(D=-L(\beta)\). Writing \(\nabla J=a\otimes I+b\otimes K\)
in the other term gives \(a(IX)=b(KX)\), hence \(b=a\circ J\).
The identity \((\nabla_{JX}J)JY=(\nabla_XJ)Y\) follows and
gives integrability of \(J\). For a Hermitian metric,
the \(W^+\) block consists of a scalar multiple of
\(\tfrac12\omega_J\otimes\omega_J-\tfrac13\operatorname{Id}\)
and off-diagonal terms between \(\omega_J\) and its
orthogonal complement. Thus \(W^+(\omega_J)=0\) makes the
entire block zero, contrary to rank two
\cite[\S4, Theorem 4.1 and proof of Proposition 4.1]{GoverNagy2007FourDimensional}.

Consequently \((d\beta)^+=0\) on \(U\).
For a nonzero smooth Bach-flat block, \(U\) is dense in \(M\):
the conformally normalized metric is real analytic in local
conformal harmonic coordinates by obstruction-flatness
\cite[Theorems 2.2 and 3.7(c)]{LassasLiimatainen2019ConformalHarmonic}.
The analytic identity theorem for a Weyl block and the conformal
invariance of its zero set exclude an open zero set on the
connected manifold. If \(\bar\beta\) is a smooth extension,
\(d\bar\beta\) is therefore anti-self-dual everywhere. Stokes
gives \(0=\int_M d\bar\beta\wedge d\bar\beta
=-\|d\bar\beta\|_2^2\). The closed extension is exact, giving
the asserted half-harmonic representative by \eqref{alt:cotton-split}.
The necessity follows by taking \(\bar\beta=df\).

In the half-harmonic scale the Bach equation is
\(C_{ikjl}P_g^{kl}=0\). On trace-free symmetric tensors this
contraction is \(W^+\) acting on the first factor in
\(S^2_0\simeq\Lambda^+\otimes\Lambda^-\).
Indeed let \(I_a,J_b\) be opposite-chirality unit quaternionic
frames with \(W^+\) diagonal in the former. The contraction on
a basis element \(I_aJ_b\) is
\(\frac12\sum_c\lambda_c I_c(I_aJ_b)I_c^{-1}\).
The conjugation fixes the term with \(c=a\) and negates the
others, yielding \(\lambda_a I_aJ_b\). Thus the trace-free
Schouten tensor is zero on the invertible locus, and then
everywhere by density. The contracted Bianchi identity gives
constant scalar curvature.
\end{proof}

The Einstein implication in the proof requires a nonzero
half-harmonic block and Bach-flatness, and applies on any connected
closed oriented four-manifold. It shows a limitation of using
stability to obtain a harmonic scale on general smooth types.
On \(N=\#_4\mathbb{CP}^2\) in the positive-signature orientation,
LeBrun supplies a positive-Yamabe self-dual conformal class
\cite[Theorem 1 and \S5, p.~235]{LeBrun1991Explicit}.
Its \(W^-=0\) makes \(E_+=12\pi^2\tau(N)=48\pi^2\)
an absolute minimum by the signature formula. It is Bach-flat
with nonnegative unrestricted Hessian. If its nonzero \(W^+\)
had a half-harmonic representative, the preceding implication
would make it positive Einstein. But then
\(4\pi^2(2\chi-3\tau)=2E_-+\int R^2/24\,dV>0\),
whereas \(\chi=6\), \(\tau=4\). On the homotopy sphere in the
subcritical band, both the vanishing residual and smooth extension
in Theorem~\ref{alt:harmonic-scale} remain hypotheses.

\subsection{The half-harmonic energy threshold}

\begin{theorem}[Strict half-harmonic bound]\label{alt:harmonic-gap}
Let \([g]\) be a smooth positive-Yamabe class on a closed
oriented homotopy four-sphere, with a smooth representative
satisfying \(\nabla^i(W^+)_{ijkl}=0\). Let \(e=E_+([g])\),
and let \(T_Y\) be the squared \(L^2\) norm of trace-free
Ricci curvature in a Yamabe minimizer of the class.
Then either the class is conformally round or
\[
 e>\frac{16\pi^2}{3}+\frac{T_Y}{6}.
\]
For the largest eigenvalue \(z\) of \(W^+\) in the nonround
case, the operator \(-6\Delta_g+R_g-6z\) has negative first
eigenvalue, and \(36\int_M z^2\,dV_g>Y([g])^2\).
\end{theorem}

\begin{proof}
In the half-harmonic representative put \(w=|W^+|_{\mathrm{HS}}\).
Use \(\Delta_g=\nabla^i\nabla_i\). The half-harmonic
Weitzenb\"ock formula and improved Kato inequality give
\[
 \tfrac12\Delta_g w^2=|\nabla W^+|^2+\tfrac12R_g w^2-18\det W^+,
 \qquad |\nabla W^+|^2\ge\tfrac53|\nabla w|^2\quad(w>0)
\]
\cite[equations (3.3), (3.7) and Corollary 3.4]{GurskyMcKeownTyrrell2026}.
Fu gives the gap with the four-index norm, whose square is
\(4w^2\) for one block
\cite[Proposition E, equations (2-1)--(2-9)]{Fu2018FourManifolds}.
For eigenvalues \(x\le y\le z\) with \(x+y+z=0\),
\[
 6zw^2-36\det W^+=12z(x-y)^2\ge0,\qquad
 6z\le 2\sqrt6\,w.
\]
The first equality uses \(w^2=2z^2-2xy\). The bound uses
\(x^2+y^2\ge z^2/2\).
Set \(u_\epsilon=w^2+\epsilon^2\),
\(f_\epsilon=u_\epsilon^{1/6}\), and \(P_z=-6\Delta_g+R_g-6z\).
Then
\[
 6\Delta_g f_\epsilon=u_\epsilon^{-5/6}
   \bigl(\Delta_g u_\epsilon-\tfrac56u_\epsilon^{-1}
                            |\nabla u_\epsilon|^2\bigr).
\]
Kato bounds the subtracted term by \(2|\nabla W^+|^2\), also
at zeroes since \(\nabla u_\epsilon=0\) there. The two preceding
displays give
\(P_z f_\epsilon\le(R_g-6z)\epsilon^2 u_\epsilon^{-5/6}
\le\|R_g\|_\infty\epsilon^{1/3}\).
Multiplication by \(f_\epsilon\) and integration bound
\(f_\epsilon\) in \(H^1\). Its uniform limit \(f=w^{1/3}\)
therefore has \(P_z\) quadratic form at most zero by lower
semicontinuity. If \(w\not\equiv0\), the Yamabe inequality
and Hölder yield
\[
 Y\|f\|_4^2
 \le \int_M(6|\nabla f|^2+R_g f^2)\,dV_g
 \le 6\int_M zf^2\,dV_g
 \le 2\sqrt6\,\|w\|_2\|f\|_4^2.
\]
In equality, Hölder forces \(w\) proportional to \(w^{2/3}\)
and hence a positive constant by continuity and connectedness.
The constant test is Yamabe minimizing, giving constant scalar
curvature \(R_g=2\sqrt6\,w\). The Weitzenb\"ock formula and
\(18\det W^+\le\sqrt6\,w^3\) force \(W^+\) to be parallel,
with eigenvalues \((-R_g/12,-R_g/12,R_g/6)\). The simple
eigenline is orientable since \(H^1(M;\mathbb Z/2)=0\).
A global unit section is a nonzero parallel self-dual form,
contrary to \(b_2(M)=0\). Thus \(e>Y^2/24\).
In the Yamabe gauge the topological identity is
\(16\pi^2=2e-T_Y/2+Y^2/24\). Combining gives the asserted
strict bound. If \(W^+=0\), the zero signature gives \(W^-=0\)
and the simply connected closed conformally flat manifold is
conformally round by the developing map.

For the spectral statement, the form tested on \(f\) gives
\(\lambda_1(P_z)\le0\). If equality holds then \(f\) is a
weak first eigenfunction. Elliptic regularity and the strong
maximum principle for the bounded Lipschitz potential give
\(f>0\). The calculation at \(\epsilon=0\) now gives
\[
 P_z f=-w^{-5/3}
     \bigl(2|\nabla W^+|^2-\tfrac{10}{3}|\nabla w|^2
                    +12z(x-y)^2\bigr).
\]
The nonnegative summands force \(x=y=-z/2\) everywhere.
A global section of norm \(\sqrt2\) of the simple eigenline
defines a compatible almost complex structure. Its first
Chern class must satisfy
\(\langle c_1^2,[M]\rangle=2\chi+3\tau=4\), impossible
when \(H^2(M;\mathbb Z)=0\). Thus the first eigenvalue is
negative. A negative-form test and Hölder give
\(Y^2<36\int z^2\,dV_g\). The sign of the operator is conformally
invariant because its quadratic form is that of the conformal
Laplacian with the weight-two potential \(6z\).
\end{proof}

Every positive Einstein metric satisfies the half-harmonic
hypothesis by the contracted Bianchi identity. Thus in the band
\(4\pi^2<E_+\le16\pi^2/3\) every Bach-flat candidate lacks
the smooth harmonic scale of Theorem~\ref{alt:harmonic-scale}.

\subsection{The positive Einstein Hessian}

\begin{proposition}[A spectral criterion for Weyl descent]
\label{alt:einstein-hessian}
On a closed oriented homotopy four-sphere let \(g\) be positive
Einstein with constant scalar \(R\). Use the positive-leading
Lichnerowicz operator on transverse-traceless tensors
\[
 \mathcal L=\nabla^*\nabla-2\mathring{\operatorname{Rm}}
                       +\tfrac R2\operatorname{Id},
 \qquad (\mathring{\operatorname{Rm}}h)_{ij}=R_{ipjq}h^{pq},
\]
where \(\mathring{\operatorname{Rm}}h=-h\) on trace-free tensors
at sectional curvature one. The Hessian of \(E_+\) on this
subspace is
\begin{equation}\label{alt:hessian}
 D^2E_+(g)[h,h]=\tfrac18
         \langle h,(\mathcal L-R/2)(\mathcal L-R/3)h\rangle_{L^2}.
\end{equation}
A negative variation modulo conformal and diffeomorphism
directions exists if and only if the transverse-traceless spectrum
of \(\mathcal L\) meets \((R/3,R/2)\). Such a direction
decreases the energy inside the positive-Yamabe classes on \(M\).
\end{proposition}

\begin{proof}
Chern--Gauss--Bonnet and signature give the functional identity
\[
 \int_M(|\operatorname{Ric}|^2-R^2/3)\,dV
 =\tfrac12\int_M |W|_{\mathrm{tens}}^2,dV-16\pi^2\chi(M)
 =4E_+-16\pi^2\chi(M).
\]
Gursky--Viaclovsky's transverse-traceless factorization for this
quadratic functional at an Einstein metric gives half the product
of the operators \(-\mathcal L+R/2\) and \(-\mathcal L+R/3\)
\cite[Theorem 3.10 and equation (3.41)]{GurskyViaclovsky2015}.
Dividing by four gives \eqref{alt:hessian}. On an eigentensor
the multiplier has negative sign exactly in the stated open
interval, and the bilinear kernel consists of the two endpoint
eigenspaces. The \(R/2\) endpoint is the infinitesimal Einstein
equation in transverse-traceless gauge.

Conformal and Lie-derivative directions are in the radical
of the Hessian by differentiating the invariances at the
critical metric. The remaining part of a smooth variation is
transverse-traceless: remove its trace and solve the orthogonal
projection onto the image of the conformal Killing operator.
Its normal operator is elliptic, and the right-hand side is
orthogonal to the conformal Killing fields. A negative
transverse-traceless tensor is realized by the path
\(g_s=g(e^{s h^\sharp}\cdot,\cdot)\), whose volume form is fixed.
For small \(|s|\), scalar curvature stays positive by compactness
and smooth dependence, giving positive Yamabe sign.
\end{proof}

Exclusion of both endpoint eigenvalues gives local Bach
rigidity, and exclusion of the full closed interval gives a
strict local minimum modulo the two gauges by the nonlinear
rigidity and local minimality results of
\cite[Theorem 1.6 and Proposition 6.4]{GurskyViaclovsky2015}.
If the open interval is empty but an endpoint occurs, the
Hessian is nonnegative and leaves those kernel directions
for higher-order analysis. The criterion concerns unrestricted
metric variations with positive Yamabe sign. It gives no path
of exact Bach-flat classes descending from an isolated value.

For any positive Einstein representative of volume \(V\),
\(R\sqrt V=Y([g])\). To see this, write
\(g=u^{-2}\widehat g\) with \(\widehat g\) a Yamabe minimizer.
The trace-free Ricci equation is
\(0=\operatorname{Ric}_{0,\widehat g}
      +2u^{-1}(\widehat\nabla^2u)_0\).
Pair with \(u\operatorname{Ric}_{0,\widehat g}\) and integrate
in \(\widehat g\). The Hessian term integrates to zero by the
contracted Bianchi identity and the constant scalar curvature.
Thus \(\widehat g\) is Einstein too. The identity
\(16\pi^2=2E_++R^2 V/24\) for the two Einstein metrics
gives the claim. In particular in the unit-volume gauge,
\(R=Y\), \(E_+=8\pi^2-Y^2/48\), and the interval test is
\(\mu/Y\in(1/3,1/2)\). For sectional curvature one,
\(R=12\) and \(\mathcal L=\nabla^*\nabla+8\), so the Hessian
is \((\nabla^*\nabla+2)(\nabla^*\nabla+4)/8\) and is
positive on this subspace. For a nonround positive Einstein
survivor the integral Weyl band gives no conclusion about an
eigenvalue in the descent interval.


\section{Outstanding problems}

The variational arguments begin with a smooth positive
subcritical conformal class on a fixed smooth type.
Existence of such a class on every homotopy four-sphere remains
open. On one fixed-type minimizing sequence with separated
quantiles, the outstanding geometric step is to extract one
physical source annulus of the width in \eqref{eq:anchored-width},
at an intermediate scale and with bounded rescaled distance to
the attaining center. Its full curvature must vanish in integral,
and its central metric must meet the first-order or critical
coordinate test of Proposition~\ref{prop:neck-critical-shell}.
The small-side estimate already determines the positions of
the two entire marks once the wide chart is supplied. Physical
transition volume and protected ball regularity give partial
tests, with the stated additional comparison hypotheses.

Selecting a minimizing sequence with vanishing integrated
Bach norm in the physical Yamabe gauge is a separate question.
For that class the quantitative descent excludes a supplied full
embedded smooth ALE limit with its specified end regularity, and
the diagonal comparison covers a specified portion of an embedded
AE end. A whole smooth marked conformal base would give another
route. Extracting such a base on weak critical curvature bounds
remains open, as does termination after a round base with a charged
puncture. In the absence of descent a nonround closed critical
class may remain.
The Einstein, harmonic-chirality and conformal criteria
above give restrictions on that alternative under their
stated hypotheses.

\paragraph{Research provenance.}
Drafted for Edward B. Baker III with OpenAI Codex
assistance on 25 November 2026 from project checkout
\texttt{5ff7014}. The exact serving model and reasoning
effort for manuscript preparation are unavailable.
The inputs are the dated neck geometry derivations and
the three joint collar assessments. The physical-shell revision
was drafted on 25 November 2026 from \texttt{483ad3c} with
OpenAI Codex assistance. It uses the incoming master E38/E41
return, the E16 and E25 local and component tests, and the
master critical-annulus and short-splice arguments. The
accompanying source inventory identifies the checked versions.
The arguments concern
smooth Riemannian sources, with quotient models and
removability used only as specified. Claims about supplied
markings, exact Bach compactness and almost Bach-flat
perturbations retain their separate hypotheses.

This revision was prepared for Edward B. Baker III with OpenAI Codex
assistance on 25 November 2026. It starts from the E59--E62 working
sources over \texttt{bfdf0660e4f7}. The E63--E67 research assessment
records the manuscript generations, independent readers and evidence
for the results used here.

\bibliographystyle{unsrt}
\bibliography{references}
\end{document}
